\exercisesection{Normed Algebras}

\begin{enumerate}
\def\labelenumi{\arabic{enumi})}
\tightlist
\item
  Let (A, e) be a unital complex algebra.
\end{enumerate}

\begin{enumerate}
\def\labelenumi{\alph{enumi})}
\tightlist
\item
  For a linear form \(\lambda \in \mathrm{A}^*\) such that \(\lambda(e) = 1\) , the following conditions are equivalent:
\end{enumerate}

\begin{enumerate}
\def\labelenumi{(\roman{enumi})}
\item
  For every \(x \in \operatorname{Ker}(\lambda)\) , one has \(x^2 \in \operatorname{Ker}(\lambda)\) ;
\item
  For every \(x \in \mathbb{A}\) , one has \(\lambda(x^{2}) = \lambda(x)^{2}\) ;
\item
  The kernel of \(\lambda\) is a right ideal of A;
\item
  The linear form \(\lambda\) is a unital morphism from A to C.
\end{enumerate}

(Observe that \(\lambda(x - \lambda(x)e) = 0\) for all \(x\in A\) .)

\begin{enumerate}
\def\labelenumi{\alph{enumi})}
\setcounter{enumi}{1}
\item
  Let \(\lambda \in \mathrm{A}^*\) a linear form such that \(\lambda(1) = 1\) and such that for every \(x \in \mathrm{A}\) , one has \(\lambda(x) \in \mathrm{Sp}_{\mathrm{A}}(x)\) . Let \(x \in \mathrm{Ker}(\lambda)\) such that \(\lambda(x^{2}) \neq 0\) The spectrum of \(x\) is not bounded in \(\mathbf{C}\) . (For \(n \geqslant 2\) integer, consider the polynomial \(p \in \mathbf{C}[\mathrm{X}]\) such that \(p(z) = \lambda((ze - x)^{n})\) for all \(z \in \mathbf{C}\) ; observe that its roots belong to the spectrum of \(x\) and bound their Euclidean norm below in terms of \(|\lambda(x^{2})|\) .)
\item
  Suppose that A is a Banach algebra. For a linear form \(\lambda\) on A to be a unital morphism from A into C, it is necessary and sufficient that \(\lambda(e)=1\) and that \(\lambda(x)\in\mathrm{Sp}_{\mathrm{A}}(x)\) for all \(x\in A\) .
\end{enumerate}

\begin{enumerate}
\def\labelenumi{\arabic{enumi})}
\setcounter{enumi}{1}
\tightlist
\item
  Let \(n \geqslant 0\) be an integer. Let \(A_n\) the commutative complex Banach algebra of functions \(f: [0,1] \to \mathbf{C}\) admitting continuous derivatives in \([0,1]\) up to order \(n\) , equipped with the norm
\end{enumerate}

\[
\| f \| = \sum_ {k = 0} ^ {n} \frac {1}{k !} \sup _ {0 \leqslant t \leqslant 1} | f ^ {(k)} (t) |
\]

(example 4 of I, p.~18). Show that \(A_{n}\) is isomorphic to \(A_{m}\) if, and only if, n = m. (However, according to Example 1 of I, p.~144, these Banach algebras all have the same character space.)

\begin{enumerate}
\def\labelenumi{\arabic{enumi})}
\setcounter{enumi}{2}
\tightlist
\item
  Let G be a group equal to its derived subgroup (A, I, p.~67). For every \(g \in G\) , we denote by \(\ell(g)\) the smallest integer \(k \geqslant 0\) such that there exist pairs \((a_{1}, b_{1}), \ldots, (a_{k}, b_{k})\) into \(G \times G\) satisfying \(g = [a_{1}, b_{1}] \cdots [a_{k}, b_{k}]\) .
\end{enumerate}

\begin{enumerate}
\def\labelenumi{\alph{enumi})}
\item
  The map \(\ell\) is a metric on G.
\item
  For every \(g \in \mathbf{G}\) , the limit
\end{enumerate}

\[
\operatorname{lsc} (g) = \lim _ {n \to + \infty} \frac {\ell (g ^ {n})}{n}
\]

exists. It is called the stable commutator length of g.

\begin{enumerate}
\def\labelenumi{\alph{enumi})}
\setcounter{enumi}{2}
\tightlist
\item
  Let S be a set and \(\mathrm{G} = \mathrm{D}(\mathrm{F}(\mathrm{S}))\) the derived subgroup of the free group generated by S. The group G is equal to its derived subgroup, and for every \(g \in G\) , \(g \neq e\) , one has \(\operatorname{lsc}(g) \geqslant 1/2\) . There exists \(g \in G\) such that \(\operatorname{lsc}(g) = 1/2\) .
\end{enumerate}

\begin{enumerate}
\def\labelenumi{\arabic{enumi})}
\setcounter{enumi}{3}
\tightlist
\item
  Let \(n \geqslant 1\) an integer. We denote \(S(n)\) the largest integer \(k \geqslant 0\) such that, for every set \(X\) of \(n\) real numbers, there exists a subset \(Y \subset X\) of cardinality \(k\) such that the equation \(x + y = z\) has no solution in \(Y\) . We therefore have \(S(n) \leqslant n\) .
\end{enumerate}

\begin{enumerate}
\def\labelenumi{\alph{enumi})}
\tightlist
\item
  Show that the limit
\end{enumerate}

\[
\sigma = \lim _ {n \to + \infty} \frac {\mathrm{S} (n)}{n}
\]

exists.

\begin{enumerate}
\def\labelenumi{\alph{enumi})}
\setcounter{enumi}{1}
\item
  Show that \(\sigma \leqslant 2 / 5\) .
\item
  Show that \(\sigma \geqslant 1 / 3\) .
\end{enumerate}

(One can show that \(\sigma = 1/3\) , cf.~S. EBERHARD, B. GREEN, F. MANNERS, Sets of integers with no large sum-free subset, Annals of Mathematics 180 (2014), 621--652.)

\begin{enumerate}
\def\labelenumi{\arabic{enumi})}
\setcounter{enumi}{4}
\tightlist
\item
  Let A be a commutative ring and \(n \geqslant 1\) an integer. For a matrix \(M = (m_{i,j})\) of type \((n, n)\) with coefficients in A, we call the permanent of M the element
\end{enumerate}

\[
\mathrm{per(M)} = \sum_ {\sigma \in \mathrm{S} _ {n}} m _ {1, \sigma (1)} \dots m _ {n, \sigma (n)}
\]

of A.

Let \(k \geqslant 1\) an integer. We denote \(X_{n,k}\) the set of matrices M of type \((n, n)\) with real coefficients such that each column of M contains k coefficients equal to 1, all the others being zero.

Show that the limit

\[
\lim _ {n \to + \infty} \left(\inf _ {\mathrm{M} \in \mathrm{X} _ {n, k}} \operatorname{per} (\mathrm{M})\right) ^ {1 / n}
\]

exists.

\begin{enumerate}
\def\labelenumi{\arabic{enumi})}
\setcounter{enumi}{5}
\tightlist
\item
  Let \(k \geqslant 1\) an integer. For every integer \(N \geqslant 1\) , we denote by \(S_k(N)\) the largest cardinality of a set \(A \subset \{1, \ldots, N\}\) that does not contain \(k\) elements in arithmetic progression, that is, for any \(n_0 \geqslant 0\) , \(h \geqslant 1\) , \(\ell \geqslant 1\) , if \(\{n_0 + h, n_0 + 2h, \ldots, n_0 + \ell h\} \subset A\) , then \(\ell < k\) .
\end{enumerate}

\begin{enumerate}
\def\labelenumi{\alph{enumi})}
\tightlist
\item
  Show that the limit
\end{enumerate}

\[
s _ {k} = \lim _ {\mathrm{N} \to + \infty} \frac {\mathrm{S} _ {k} (\mathrm{N})}{\mathrm{N}}
\]

exists.

\begin{enumerate}
\def\labelenumi{\alph{enumi})}
\setcounter{enumi}{1}
\tightlist
\item
  Show that \(s_{1}=s_{2}=0\) One can show (“Szemerédi's theorem”, cf.~E. SZEMERÉDI, On sets of integers containing no k elements in arithmetic progression, Acta Arithmetica XXVII, 1975, 199--245) that \(s_{k}=0\) for all \(k\geqslant1\) ; for the case k=3, cf.~exercise 23 of II, p.~270.
\end{enumerate}

\begin{enumerate}
\def\labelenumi{\arabic{enumi})}
\setcounter{enumi}{6}
\tightlist
\item
  Let A be a normed algebra.
\end{enumerate}

\begin{enumerate}
\def\labelenumi{\alph{enumi})}
\tightlist
\item
  Let S be a nonempty bounded subset of A. Show that the limit
\end{enumerate}

\[
\varrho (\mathrm{S}) = \lim _ {n \to + \infty} \sup _ {x \in \mathrm{S} ^ {n}} \| x \| ^ {1 / n}
\]

exists and that \(\varrho(\mathrm{S}) = \varrho(x)\) if S is reduced to a single element x. We say that \(\varrho(\mathrm{S})\) is the spectral radius of the part S.

\begin{enumerate}
\def\labelenumi{\alph{enumi})}
\setcounter{enumi}{1}
\tightlist
\item
  Let N be the set of norms p on A, equivalent to the norm of A, and such that A, equipped with p, is a normed algebra. Let X be a subset of A. For there to exist \(p \in N\) such that \(p(x) \leqslant 1\) for all \(x \in X\) , it is necessary and sufficient that there exist a real number \(c \geqslant 0\) such that
\end{enumerate}

\[
\sup _ {n \geqslant 1} \sup _ {x \in X ^ {n}} \| x \| \leqslant c.
\]

\begin{enumerate}
\def\labelenumi{\alph{enumi})}
\setcounter{enumi}{2}
\tightlist
\item
  Show that
\end{enumerate}

\[
\varrho (\mathrm{S}) = \inf _ {p \in \mathrm{N}} \sup _ {x \in \mathrm{S}} p (x).
\]

\begin{enumerate}
\def\labelenumi{\arabic{enumi})}
\setcounter{enumi}{7}
\item
  Let A be the algebra of continuous complex-valued functions on \(\mathbf{R}\) tending to 0 at infinity, equipped with the norm \(\|f\| = \sup_{t\in\mathbf{R}} |f(t)|\) . Then \(\widetilde{\mathrm{A}}\) is identified with the algebra of continuous complex-valued functions on \(\mathbf{R}\) tending toward a finite limit at infinity. On \(\widetilde{\mathrm{A}}\) , the norm \(\|g\| = \sup_{t\in\mathbf{R}} |g(t)|\) differs from the norm defined in no. 1 of I, p.~15.
\item
  Let A be a normed algebra and \(b = \inf_{x \neq 0} \|x^{2}\|/\|x\|^{2}\) . Then
\end{enumerate}

\[
b \leqslant \inf _ {x \neq 0} \frac {\varrho (x)}{\| x \|} \leqslant \sqrt {b}.
\]

\begin{enumerate}
\def\labelenumi{\arabic{enumi})}
\setcounter{enumi}{9}
\item
  Let A be a normed algebra and \(x, y \in A\) . We have \(\varrho(xy) = \varrho(yx)\) .
\item
  Let H be a Hilbert space of countable type and \((f_{n})_{n\geqslant1}\) an orthonormal basis of H. Let A be the normed algebra \(\mathcal{L}(\mathrm{H})\) . Let \(x\in A\) the element defined by \(x(f_{i})=2^{-i}f_{i+1}\) . Show that x is quasi-nilpotent, but not nilpotent.
\end{enumerate}

\begin{enumerate}
\def\labelenumi{¶ \arabic{enumi})}
\setcounter{enumi}{11}
\tightlist
\item
  Let H be a separable Hilbert space and let \((f_n)_{n \geqslant 1}\) be an orthonormal basis of H. One defines the numbers \((\alpha_m)_{m \geqslant 1}\) by \(\alpha_m = e^{-k}\) if \(m\) is the product of \(2^k\) by an odd number.
\end{enumerate}

\begin{enumerate}
\def\labelenumi{\alph{enumi})}
\tightlist
\item
  Let \(x \in \mathcal{L}(\mathrm{H})\) defined by \(x(f_m) = \alpha_m f_{m+1}\) for \(m \geqslant 1\) Prove that \(x\) is not quasi-nilpotent. (Observe that
\end{enumerate}

\[
\left\| x ^ {n} \right\| = \sup _ {m} \bigl (\alpha_ {m} \alpha_ {m + 1} \dots \alpha_ {m + n - 1} \bigr),
\]

and evaluate \(\alpha_{1}\alpha_{2}\ldots\alpha_{2^{t}-1}\) .)

\begin{enumerate}
\def\labelenumi{\alph{enumi})}
\setcounter{enumi}{1}
\tightlist
\item
  Let \(x_{k} \in \mathcal{L}(\mathrm{H})\) defined by \(x_{k}(f_{m}) = 0\) if m is the product of \(2^{k}\) by an odd number, and \(x_{k}(f_{m}) = \alpha_{m} f_{m+1}\) otherwise. Show that \(x_{k}\) is nilpotent and that \((x_{k})\) converges to x as k tends to \(+\infty\) .
\end{enumerate}

\begin{enumerate}
\def\labelenumi{\arabic{enumi})}
\setcounter{enumi}{12}
\item
  Let A be a unital Banach algebra. Let \(x \in A\) such that \(\|x - 1\| < 1\) . There exists \(y \in A\) such that \(y^{2} = x\) .
\item
  Let A be a unital complex normed algebra. If \(\|x^{-1}\| = \|x\|^{-1}\) for every invertible element x of A, one has A = C · 1. (Reduce to the case where A is complete. Let \(x \in A\) , at a distance \(\alpha > 0\) from C · 1. If \(\lambda_{0} \in C\) is such that \(x - \lambda_{0}\) were invertible, then \(x - \lambda\) would be invertible for \(|\lambda - \lambda_{0}| < \alpha\) . Deduce that the spectrum of x would be empty.)
\item
  Let A be a unital complex normed algebra in which the only left topological divisor of zero is 0 and the only right topological divisor of zero is 0. Then \(A = C \cdot 1\) . (If \(x \in A\) , and if \(\lambda\) belongs to the boundary of \(\mathrm{Sp}_{\widehat{\mathrm{A}}}(\boldsymbol{x})\) , then \(x - \lambda\) is a topological divisor of zero in \(\widehat{A}\) , hence in A, therefore \(x = \lambda\) .)
\item
  Let A be a unital complex Banach algebra. Let \(a\) an element of A. The singular spectrum of A is called the set of \(\lambda \in \mathbf{C}\) such that \(a - \lambda e\) is a topological zero divisor on the left and on the right in A.
\end{enumerate}

\begin{enumerate}
\def\labelenumi{\alph{enumi})}
\item
  Show that the singular spectrum of A is contained in the spectrum of A, and contains the boundary of the spectrum of A.
\item
  Let \(f \in C[X]\) a polynomial. The image under f of the singular spectrum of a is the singular spectrum of \(f(a)\) .
\end{enumerate}

\begin{enumerate}
\def\labelenumi{\arabic{enumi})}
\setcounter{enumi}{16}
\tightlist
\item
  Let A be a normed algebra. For every \(x \in A\) , we set:
\end{enumerate}

\[
\lambda (x) = \inf _ {y \neq 0} \frac {\| x y \|}{\| y \|} \quad \lambda^ {\prime} (x) = \inf _ {y \neq 0} \frac {\| y x \|}{\| y \|}.
\]

\begin{enumerate}
\def\labelenumi{\alph{enumi})}
\tightlist
\item
  We have
\end{enumerate}

\[
| \lambda (x) - \lambda (y) | \leqslant \| x - y \|, \quad | \lambda^ {\prime} (x) - \lambda^ {\prime} (y) | \leqslant \| x - y \|,
\]

\[
\lambda (x) \lambda (y) \leqslant \lambda (x y) \leqslant \| x \| \lambda (y), \quad \lambda^ {\prime} (x) \lambda^ {\prime} (y) \leqslant \lambda^ {\prime} (x y) \leqslant \lambda^ {\prime} (x) \| y \|.
\]

\begin{enumerate}
\def\labelenumi{\alph{enumi})}
\setcounter{enumi}{1}
\tightlist
\item
  The set of left (resp. right) topological zero divisors in A is closed.
\end{enumerate}

\begin{enumerate}
\def\labelenumi{\arabic{enumi})}
\setcounter{enumi}{17}
\item
  Let A be a unital Banach algebra. The set of \(x\) which are neither invertible nor topological zero divisors, is open in A.
\item
  Let X be a complex Banach space and \(x \in A = \mathcal{L}(X)\) .
\end{enumerate}

\begin{enumerate}
\def\labelenumi{\alph{enumi})}
\tightlist
\item
  If x is not invertible, then x is a left or right topological zero divisor. (Consider successively the following cases:)
\end{enumerate}

\(1^{\circ}\) ) \(x\) is not injective;

\(2^{\circ})\overline{x(\mathrm{X})}\neq \mathrm{X};\)

\(3^{\circ}) x \text{ is injective, } x(\mathrm{X}) \text{ is dense in X and distinct from X.}\)

\begin{enumerate}
\def\labelenumi{\alph{enumi})}
\setcounter{enumi}{1}
\tightlist
\item
  If x maps X bicontinuously onto a closed vector subspace of X distinct from X, or if x is non-injective with image X, then x is interior to the set of non-invertible elements of A.
\end{enumerate}

\begin{enumerate}
\def\labelenumi{\arabic{enumi})}
\setcounter{enumi}{19}
\item
  In the algebra A of Example 9 of I, p.~20, the identity function z is neither invertible nor a topological zero divisor.
\item
  Let A be a unital normed algebra. Let B be the normed algebra \(\mathcal{L}(\mathrm{A})\) . Then the image of A under the morphism \(x \mapsto \gamma_x\) is a full subalgebra of B. Therefore, if \(x \in \mathrm{A}\) , one has \(\mathrm{Sp}_{\mathrm{A}}(x) = \mathrm{Sp}_{\mathrm{B}}(\gamma_x)\) .
\end{enumerate}

\begin{enumerate}
\def\labelenumi{¶ \arabic{enumi})}
\setcounter{enumi}{21}
\tightlist
\item
  Let A be a Banach algebra.
\end{enumerate}

\begin{enumerate}
\def\labelenumi{\alph{enumi})}
\tightlist
\item
  Let S be the set of bounded sequences of elements of A, with the norm \(\|(\boldsymbol{x}_{n})\| = \sup \|x_{n}\|\) . Let R be the set of \((x_{n}) \in \mathrm{S}\) such that \(\|x_{n}\|\) tends to 0. Then S is a Banach algebra and R is a closed two-sided ideal of S.
\end{enumerate}

Let \(S' = \text{S/R}\) and \(\varphi: \text{S} \to \text{S}'\) the canonical homomorphism.

\begin{enumerate}
\def\labelenumi{\alph{enumi})}
\setcounter{enumi}{1}
\item
  The map \(\theta\) from A into S' defined by \(\theta(x) = \varphi(x, x, x, \ldots)\) is an isometric isomorphism of A onto a subalgebra of S'. Every left topological divisor of zero in S' is a left divisor of zero. For \(x \in A\) to be a left topological divisor of zero, it is necessary and sufficient that \(\theta(x)\) be a left divisor of zero in S'.
\item
  Assume that A has a unit element. Prove that there exists a Banach algebra B and an isometric isomorphism of A onto a subalgebra of B such that an element of A is non-invertible if and only if its image in B is a left or right divisor of zero. (Use (a) and Exercises 19 and 21.)
\end{enumerate}

\begin{enumerate}
\def\labelenumi{\arabic{enumi})}
\setcounter{enumi}{22}
\tightlist
\item
  Let A be a normed algebra and R its radical.
\end{enumerate}

\begin{enumerate}
\def\labelenumi{\alph{enumi})}
\item
  If \(x \in R\) , then x is quasi-nilpotent. (We have \(\mathrm{Sp}_{\mathrm{A}}^{\prime}(x) = \{0\}\) and a fortiori \(\mathrm{Sp}_{\widehat{\mathrm{A}}}^{\prime}(x) = \{0\}\) .)
\item
  Suppose that A is a Banach algebra. Let I be a left ideal of A all of whose elements are quasi-nilpotent. Then \(I \subset R\) (We may assume that A has a unit element. If \(x \in R\) and \(a \in A\) , one has \(\mathrm{Sp}_{\mathrm{A}}(ax) = \{0\}\) , hence 1 - ax is invertible; hence \(x \in R\) .
\end{enumerate}

\begin{enumerate}
\def\labelenumi{\arabic{enumi})}
\setcounter{enumi}{23}
\item
  Let A be a complex Banach algebra, B a complex algebra without radical, and \(\varphi\) a morphism from A onto B. The kernel of \(\varphi\) is closed.
\item
  Let A be a complex Banach algebra and \(\pi\) a nonzero irreducible representation of A in a complex vector space X. Let \(\xi_{0}\) be a nonzero element of X.
\end{enumerate}

\begin{enumerate}
\def\labelenumi{\alph{enumi})}
\item
  The annihilator I of \(\xi_0\) is a maximal regular left ideal of A; it is closed.
\item
  The map \(x \mapsto x\xi_{0}\) defines, by passage to the quotient, an isomorphism \(\varphi\) of the A-module A/I onto the A-module X.
\item
  If one transports by \(\varphi\) the norm of the Banach space A/I, then X becomes a Banach space, and \(\| \pi (x)\| \leqslant \| x\|\) for all \(x\in \mathrm{A}\) .
\item
  If A is primitive (that is, \(\{0\}\) is a primitive ideal, cf. definition I, p. 11, cf. also A, VIII, p. 433, exerc. 6), then \(A = \{0\}\) or \(C \cdot 1\) according as A does or does not have a unit element. (Use the preceding and cor. 4 of I, p. 26).
\end{enumerate}

\begin{enumerate}
\def\labelenumi{¶ \arabic{enumi})}
\setcounter{enumi}{25}
\tightlist
\item
  Let A be a Banach algebra and R its radical.
\end{enumerate}

\begin{enumerate}
\def\labelenumi{\alph{enumi})}
\tightlist
\item
  Prove that if \(r \in R\) , the series
\end{enumerate}

\[
- \frac {1}{2} \sum_ {k = 1} ^ {\infty} \binom{1 / 2}{k} (- 4 r) ^ {k}
\]

converges to an element \(x \in R\) such that \(x^{2} - x + r = 0\) .

\begin{enumerate}
\def\labelenumi{\alph{enumi})}
\setcounter{enumi}{1}
\tightlist
\item
  Let u be an element of A whose class in A/R is idempotent. There exists an idempotent of A congruent to u modulo R. (We may assume that A has a unit element. Let \(q = u - u^{2} \in R\) and \(x \in R\) a solution, constructed with the aid of a), of \(x^{2} - x - q(1 - 4q)^{-1} = 0\) . Then \(u - (2u - 1)x\) answers the question.)
\end{enumerate}

\begin{enumerate}
\def\labelenumi{\arabic{enumi})}
\setcounter{enumi}{26}
\item
  Let A be a complex unital Banach algebra, \(x \in A\) and \(\zeta\) a boundary point of \(\mathrm{Sp}_{\mathrm{A}}(x)\) . The resolvent of x cannot be extended to a continuous function at \(\zeta\) .
\item
  Let A be a complex unital Banach algebra, \(\Delta\) an open subset of C and \(\lambda \mapsto \mathrm{R}(\lambda)\) a map of \(\Delta\) to A such that
\end{enumerate}

\[
\mathrm{R} (\lambda) - \mathrm{R} (\mu) = - (\lambda - \mu) \mathrm{R} (\lambda) \mathrm{R} (\mu)\tag{5}
\]

for all \(\lambda,\mu\in\Delta\) .

\begin{enumerate}
\def\labelenumi{\alph{enumi})}
\tightlist
\item
  Let \(\lambda_0\in \Delta\) . For every \(\lambda \in \Delta\) such that \(|\lambda -\lambda_0|\cdot \| \mathrm{R}(\lambda)\| <  1\) , we have:
\end{enumerate}

\[
\mathrm{R} (\lambda) = \sum_ {n = 0} ^ {\infty} (\lambda_ {0} - \lambda) ^ {n} \mathrm{R} (\lambda_ {0}) ^ {n + 1}.
\]

\begin{enumerate}
\def\labelenumi{\alph{enumi})}
\setcounter{enumi}{1}
\tightlist
\item
  For every \(\lambda \in \Delta\) and every integer \(n\geqslant 0\) , one has
\end{enumerate}

\[
\frac {\partial^ {n}}{\partial \lambda^ {n}} \mathrm{R} (\lambda) = (- 1) ^ {n} n! \mathrm{R} (\lambda) ^ {n + 1}.
\]

\begin{enumerate}
\def\labelenumi{\alph{enumi})}
\setcounter{enumi}{2}
\tightlist
\item
  If R is holomorphic at infinity, then there exists \(z, j, x \in \mathrm{A}\) such that \(z^2 = 0\) , \(j^2 = j\) , \(zj = jz = 0\) , \(x \in j\mathrm{Aj}\) , and \(\mathrm{R}(\lambda) = z + j\mathrm{R}(\lambda, x)\) for \(|\lambda|\) sufficiently large. (Write \(\mathrm{R}(\lambda) = \sum_{n=0}^{\infty} c_n \lambda^{-n}\) for \(|\lambda| > \lambda_0\) , and express that R satisfies (5). One may set \(c_0 = z\) , \(c_1 = j\) , \(c_2 = x\) .)
\end{enumerate}

\begin{enumerate}
\def\labelenumi{\alph{enumi})}
\setcounter{enumi}{3}
\tightlist
\item
  For there to exist \(a \in A\) such that \(R\) is the restriction to \(\Delta\) of the resolvent \(\lambda \mapsto R(a, \lambda)\) of \(a\) , it is necessary and sufficient that there exist \(\lambda_0 \in \Delta\) such that \(R(\lambda_0)\) be invertible.
\end{enumerate}

\begin{enumerate}
\def\labelenumi{\alph{enumi})}
\setcounter{enumi}{4}
\tightlist
\item
  For every \(x \in A\) , the function \(\lambda \mapsto \sum_{n=0}^{\infty} (\lambda_0 - \lambda)^n x^{n+1}\) satisfies (5) for \(|\lambda - \lambda_0| \cdot \|x\| < 1\) .
\end{enumerate}

\begin{enumerate}
\def\labelenumi{\arabic{enumi})}
\setcounter{enumi}{28}
\tightlist
\item
  \begin{enumerate}
  \def\labelenumii{\alph{enumii})}
  \tightlist
  \item
    Let E be a complex Banach space, \(f: \mathbf{C} \to \mathbf{E}\) an entire function and \(\theta \mapsto \mathrm{P}(\theta) = \sum_{\mu=k}^{m} c_{\mu} e^{\mu i\theta}\) a trigonometric polynomial ( \(c_k \neq 0\) ). We assume that \(|\mathrm{P}(\theta)| \cdot \|f(re^{i\theta})\| \leqslant \mathrm{Mr}^{\alpha}\) for \(r \geqslant r_0\) (M, \(r_0\) , \(\alpha\) being constants \(\geqslant 0\) ). Then \(f\) is a polynomial of degree \(\leqslant \alpha\) . (Let \(f(\zeta) = \sum_{\nu=0}^{\infty} a_{\nu} \zeta^{\nu}\) , and
  \end{enumerate}
\end{enumerate}

\[
\mathrm{J} (r, p) = r ^ {- p} \int_ {0} ^ {2 \pi} \mathrm{P} (\theta) f (r e ^ {i \theta}) e ^ {- (k + p) i \theta} d \theta .
\]

Compute the limit of \(\mathrm{J}(r,p)\) when \(r\to+\infty\) in two different ways for \(p>\alpha.\)

\begin{enumerate}
\def\labelenumi{\alph{enumi})}
\setcounter{enumi}{1}
\tightlist
\item
  With the notations of \(a\) ), if \(\alpha\) is an integer and if, for every \(\theta \in \mathbf{R}\) , one has
\end{enumerate}

\[
\lim _ {r \to + \infty} r ^ {- \alpha} | \mathrm{P} (\theta) | \| f (r e ^ {i \theta}) = 0,
\]

then \(f\) is a polynomial of degree \(< \alpha\) .

\begin{enumerate}
\def\labelenumi{\alph{enumi})}
\setcounter{enumi}{2}
\tightlist
\item
  Let A be a unital complex Banach algebra and x a quasi-nilpotent element of A. We assume that, for some integer n \textgreater{} 0, and for every \(\theta \in R\) , one has
\end{enumerate}

\[
\lim _ {r \to + \infty} r ^ {- n} | \cos^ {n} \theta | \cdot \| (1 - r e ^ {i \theta} x) ^ {- 1} \| = 0.
\]

Then \(x^n = 0\) . (Use \(b\) ).)

\begin{enumerate}
\def\labelenumi{¶ \arabic{enumi})}
\setcounter{enumi}{29}
\tightlist
\item
  a) Let E be a complex Banach space and \(x: \mathbf{C}=\{1\}\to\mathrm{E}\) an entire function of \(1/(\zeta-1)\) . One denotes by \(x(\zeta)=\sum_{n=0}^{\infty}a_{n}\zeta^{n}\) , \(x(\zeta)=\sum_{n=0}^{\infty}b_{n}\zeta^{-n}\) the expansions of \(x\) for \(|\zeta|<1\) , \(|\zeta|>1\) respectively. We assume that there exists \(\alpha>0\) such that \(\|a_{n}\|=o(n^{\alpha})\) , \(\|b_{n}\|=o(n^{\alpha})\) . Then \(x(\zeta)\) is a polynomial in \(1 / (\zeta - 1)\) of degree \(< \alpha + 1\) . (Let \(\varepsilon > 0\) . We have \(\|x(\zeta)\| \leqslant \varepsilon(1 - r)^{-1 - \alpha}\) for \(r_{\varepsilon} \leqslant r = |\zeta| < 1\) , \(\|x(\zeta)\| \leqslant \varepsilon(1 - r^{-1})^{-1 - \alpha}\) for \(1 < r \leqslant r_{\varepsilon}'\) . Setting \(1 - \zeta = (\omega + \frac{1}{2})^{-1}\) , \(\omega = \mathrm{Re}^{i\theta}\) , show that, for \(0 \leqslant r < 1\) and \(\mathrm{R} \geqslant 1\) , one has
\end{enumerate}

\[
(1 - r) ^ {- 1} \leqslant \frac {9}{4} \frac {\mathrm{R}}{\cos \theta}
\]

and that, for 1 \textless{} r and R ≥ 1,

\[
(1 - r ^ {- 1}) ^ {- 1} \leqslant \frac {9}{4} \frac {\mathrm{R}}{\cos \theta}.
\]

Set \(y(\omega) = x(\zeta)\) and show that \(R^{-\alpha - 1}|\cos \theta|^{\alpha + 1}\| y(Re^{i\theta})\|\) tends to 0 as \(R\) tends to \(+\infty\) , uniformly in \(\theta\) . Then apply exerc. 29, b).

Let A be a complex unital Banach algebra, \(q \in A\) a quasi-nilpotent element and \(x = 1 + q\) .

\begin{enumerate}
\def\labelenumi{\alph{enumi})}
\setcounter{enumi}{1}
\item
  For \(q^{N}=0\) , it is necessary and sufficient that \(\|x^{\pm n}\|=o(n^{N})\) when n tends to \(+\infty\) . (To see that the condition is sufficient, show, by applying a), that \(\mathrm{R}(\lambda,x)\) is a polynomial in \(1/(\lambda-1)\) of degree \(\leqslant N\) . On the other hand, \(\mathrm{R}(\lambda,x)=\sum_{n=0}^{\infty}q^{n}(\lambda-1)^{-n-1}\) .)
\item
  Let R be the radical of A and G a subgroup of the group of invertible elements of A. If, for every \(x \in G\) , one has \(\|x^{n}\| = o(n)\) , \(\|x^{-n}\| = o(n)\) for \(n \to +\infty\) , then the restriction to G of the canonical map \(A \to A/R\) is injective.
\end{enumerate}

\begin{enumerate}
\def\labelenumi{\arabic{enumi})}
\setcounter{enumi}{30}
\tightlist
\item
  \begin{enumerate}
  \def\labelenumii{\alph{enumii})}
  \tightlist
  \item
    Let A be a normed algebra, and x, y two elements of A such that xy - yx = y. Prove that \(y^{n} = 0\) for \(n > 2\|x\|\) . (Use the formula \(xy^{n} - y^{n}x = ny^{n}\) ).
  \end{enumerate}
\end{enumerate}

\begin{enumerate}
\def\labelenumi{\alph{enumi})}
\setcounter{enumi}{1}
\item
  Let L be a finite-dimensional Lie algebra over C. Prove that, if L is not nilpotent, it contains two nonzero elements x, y such that \([x, y] = y\) . (Use the fact that there exists an \(x \in L\) such that \(\operatorname{ad}(x)\) is not nilpotent.)
\item
  Let L be a real or complex finite-dimensional Lie algebra such that the enveloping algebra of L possesses a norm compatible with its algebra structure. Prove that L is nilpotent. (Use a) and b).
\end{enumerate}

\begin{enumerate}
\def\labelenumi{\alph{enumi})}
\setcounter{enumi}{3}
\tightlist
\item
  Let V be a finite-dimensional real or complex vector space, \(\mathfrak{L}\) a Lie subalgebra of \(\mathfrak{gl}(\mathrm{V})\) consisting of nilpotent endomorphisms, and \(x_{1},\ldots ,x_{n}\) elements of \(\mathfrak{L}\) . Let \(\alpha >0\) . Prove that there exists a Hilbert space structure on V such that \(\| x_i\| \leqslant \alpha\) for \(1\leqslant i\leqslant n\) .
\end{enumerate}

\begin{enumerate}
\def\labelenumi{\alph{enumi})}
\setcounter{enumi}{4}
\tightlist
\item
  Let \(\mathfrak{L}\) be a real or complex finite-dimensional nilpotent Lie algebra, and U its enveloping algebra. Prove that there exists on U a norm compatible with its algebra structure. (Using the method of LIE, I, §3, exerc. 5 and §7, exerc. 3, show that there exists a family \((\pi_{\lambda})\) of representations of \(\mathfrak{L}\) in finite-dimensional spaces \(V_{\lambda}\) , having the following property:
\end{enumerate}

for all \(u \in U\) , there exists an \(\lambda\) such that \(\pi_{\lambda}(u) \neq 0\) , and \(\pi_{\lambda}(\mathfrak{L})\) consists, for every \(\lambda\) , of nilpotent endomorphisms. By endowing each \(V_{\lambda}\) with a Hilbert structure furnished by d), one obtains in the Hilbert space V, the Hilbert sum of the \(V_{\lambda}\) , a representation \(\pi\) of L by continuous endomorphisms, whence an injective morphism \(\varphi\) of U into \(\mathcal{L}(V)\) . Set \(\|u\|_{U} = \|\varphi(u)\|\) for all \(u \in U\) ).

\begin{enumerate}
\def\labelenumi{\alph{enumi})}
\setcounter{enumi}{5}
\tightlist
\item
  Let A be a complex unital normed algebra and x and y elements of A. If \(A \neq \{0\}\) , one has xy - yx ≠ 1. (First proof: one has \(\mathrm{Sp}_{\mathrm{A}}^{\prime}(xy) = \mathrm{Sp}_{\mathrm{A}}^{\prime}(yx)\) ; if xy = yx + 1, then \(\mathrm{Sp}_{\mathrm{A}}(xy)\) is deduced from \(\mathrm{Sp}_{\mathrm{A}}(yx)\) by the translation \(z \mapsto z + 1\) ; deduce that \(\mathrm{Sp}_{\mathrm{A}}(xy)\) is unbounded, which is absurd. Second proof: if \([x, y] = 1\) , one has \([xy, y] = y\) , whence \(y^{n} = 0\) for n sufficiently large by a); since \([x, y^{p}] = py^{p-1}\) deduce that y = 0).
\end{enumerate}

\begin{enumerate}
\def\labelenumi{\arabic{enumi})}
\setcounter{enumi}{31}
\item
  Let A be a normed algebra and \(x \in A\) . One says that x is topologically nilpotent if \(x^{n}\) tends to 0 when n tends to \(+\infty\) . Show that x is topologically nilpotent if and only if \(\varrho(x) < 1\) .
\item
  Let A be a commutative unital Banach algebra over C. Suppose that every closed ideal I of A is finitely generated (that is, there exist \(n \in N\) and \(x_{1}, \ldots, x_{n} \in A\) such that \(I = Ax_{1} + \cdots + Ax_{n}\) ).
\end{enumerate}

\begin{enumerate}
\def\labelenumi{\alph{enumi})}
\tightlist
\item
  Every ideal I of A is closed. (Write \(\bar{\mathrm{I}} = \mathrm{A}x_{1} + \dots +\mathrm{A}x_{n}\) with \(x_{1},\ldots ,x_{n}\) in A. Let \((x_{ip})_{p\geqslant 1}\) be a sequence of elements of I tending to \(x_{i}\) . For every \((y_{1},\ldots ,y_{n})\in \mathrm{A}^{n}\) and every \(p\geqslant 1\) , set
\end{enumerate}

\[
\begin{array}{c} \varphi (y _ {1}, \ldots , y _ {n}) = x _ {1} y _ {1} + \dots + x _ {n} y _ {n} \in \bar {\mathrm{I}} \\ \varphi_ {p} (y _ {1}, \ldots , y _ {n}) = x _ {i p} y _ {1} + \dots + x _ {n p} y _ {n} \in \mathrm{I}. \end{array}
\]

We have \(\varphi\in\mathcal{L}(A^{n},\bar{I})\) , \(\varphi_{p}\in\mathcal{L}(A^{n},\bar{I})\) , \(\|\varphi-\varphi_{p}\|\) tends to 0 as \(p\) tends to \(+\infty\) , and \(\varphi\) is surjective. Deduce that \(\varphi_{p}\) is surjective for \(p\) large enough, considering the transposes of \(\varphi\) and \(\varphi_{p}\) Hence I= \(\bar{I}\) .

\begin{enumerate}
\def\labelenumi{\alph{enumi})}
\setcounter{enumi}{1}
\item
  If A is an integral domain, \(A = C \cdot 1\) (Use a) and Exerc. 15.)
\item
  If the only nilpotent element of A is 0, there exists an integer \(n \geqslant 0\) such that A is isomorphic to \(\mathbf{C}^n\) . (Since A is Noetherian, \(\{0\}\) is the intersection of prime ideals \(\mathfrak{p}_1, \ldots, \mathfrak{p}_m\) ; these are closed by \(a\) ), and the algebras \(\mathrm{A}/\mathfrak{p}_i\) are isomorphic to \(\mathbf{C}\) according to \(b\) .)
\end{enumerate}

\begin{enumerate}
\def\labelenumi{\alph{enumi})}
\setcounter{enumi}{3}
\tightlist
\item
  We have \(\dim_{\mathbf{C}}\mathrm{A} < +\infty\) . (Let N be the set of nilpotent elements of A, which is a (closed) ideal of A. Then \(\dim (\mathrm{A / N}) < +\infty\) according to \(c\) ). Since N is a finitely generated ideal, \(\mathrm{N}^n = 0\) for \(n\) sufficiently large. Finally, each \(\mathrm{N}^i /\mathrm{N}^{i + 1}\) is a finitely generated module over A/N, hence \(\dim (\mathrm{N}^i /\mathrm{N}^{i + 1}) < +\infty\) .)
\end{enumerate}

\begin{enumerate}
\def\labelenumi{\arabic{enumi})}
\setcounter{enumi}{33}
\tightlist
\item
  Let A be a unital algebra over C, equipped with a separated locally convex topology such that multiplication in A is separately continuous.
\end{enumerate}

. An element \(a \in A\) is said to be regular if there exists \(r \geqslant 0\) such that \(a - \lambda\) let be invertible for \(|\lambda| \geqslant r\) , and such that the set of \((a - \lambda)^{-1}\) , for \(|\lambda| \geqslant r\) , be bounded.

\begin{enumerate}
\def\labelenumi{\alph{enumi})}
\item
  If a is regular, \(\lambda(a-\lambda)^{-1}\) remains bounded for \(|\lambda|\geqslant r\) .
\item
  If the mapping \(x \mapsto x^{-1}\) is defined in a neighborhood of 1 and continuous at the point 1, every element of A is regular.
\item
  For every \(a \in A\) , let \(U_{a}\) the set of \(\lambda \in C\) such that \((a - \mu)^{-1}\) exists and is bounded for all \(\mu\) of a neighborhood of \(\lambda\) . Let \(S_{a} = C - U_{a}\) . Then \(S_{a}\) is closed; and, if a is regular, \(S_{a}\) is compact nonempty (provided that \(A \neq \{0\}\) ).
\item
  Let \(a \in A\) . The function \(\lambda \mapsto (a - \lambda)^{-1}\) is holomorphic in \(U_{a}\) .
\item
  Let \(a \in A\) . For \(\lambda \in U_{a}\) , it is necessary and sufficient that \(a - \lambda\) has a regular inverse.
\item
  If every element of A is regular, and if A is a field, then we have \(A = C \cdot 1\) .
\end{enumerate}

\begin{enumerate}
\def\labelenumi{\arabic{enumi})}
\setcounter{enumi}{34}
\tightlist
\item
  Let A be an algebra over R which is a field, equipped with a separated locally convex topology such that multiplication in A is continuous, and such that the map \(x \mapsto x^{-1}\) be continuous at 1. Then A is R-isomorphic either to R, to C, or to the quaternion field H. (If A is commutative and there exists \(u \in A\) with \(u^{2} = -1\) give A a structure of algebra over C, and apply exerc. 34, f).) If A is commutative and there exists no element \(u \in A\) such that \(u^{2} = -1\) apply exerc. 34, f) to the field \(A \otimes_{R} C\) If A is noncommutative, argue as in AC, VI, §6, no. 4, th. 1, third case.)
\end{enumerate}

\begin{enumerate}
\def\labelenumi{¶ \arabic{enumi})}
\setcounter{enumi}{35}
\tightlist
\item
  Let A be the algebra over C consisting of the restrictions to {[}0,1{]} of rational functions in one variable with complex coefficients. This algebra is a field.
\end{enumerate}

\begin{enumerate}
\def\labelenumi{\alph{enumi})}
\item
  We endow A with the topology of convergence in measure (INT, IV, §5, no. 11). Then the mappings \((x,y)\mapsto x+y\) and \((x,y)\mapsto xy\) of A × A into A are continuous, and the map \(x\mapsto x^{-1}\) from A* to A is continuous. The topology of A is not locally convex.
\item
  For \(n \in N^{*}\) , we define the sequence \((w_{r,n})_{r \in \mathbf{Z}}\) by \(w_{r,n} = (-r + 1)^{n(-r+1)}\) if r \textless{} 0, \(w_{0,n} = 1\) and \(w_{r,n} = (r + 1)^{-(r+1)/n}\) if r \textgreater{} 0. If \(f \in A\) , set \(p_n(f) = \sum_r w_{r,n} |a_r| < +\infty\) , where \(\sum a_r t^r\) is the Laurent expansion of \(f(t)\) at 0. Then the \(p_n\) are semi-norms that define on A a metrizable locally convex topology. The map \((x, y) \mapsto xy\) The mapping from A × A to A is continuous. The map \(x \mapsto x^{-1}\) The map from A* into A is not continuous. The algebra A is not complete for this topology.
\end{enumerate}

\begin{enumerate}
\def\labelenumi{\arabic{enumi})}
\setcounter{enumi}{36}
\tightlist
\item
  \begin{enumerate}
  \def\labelenumii{\alph{enumii})}
  \tightlist
  \item
    Let A be an algebra over C equipped with a locally convex topology. The following conditions are equivalent:
  \end{enumerate}
\end{enumerate}

\begin{enumerate}
\def\labelenumi{(\roman{enumi})}
\item
  there exists a fundamental system \((\mathrm{U}_{i})\) of convex balanced neighborhoods of 0 such that \(U_{i}U_{i} \subset U_{i}\) for all i ;
\item
  A is isomorphic to a subalgebra of a product of normed algebras.
\item
  the topology of A can be defined by a family of semi-norms \(p_{i}\) satisfying \(p_{i}(xy) \leqslant p_{i}(x)p_{i}(y)\) for all \(x, y \in A\) .
\end{enumerate}

If A satisfies these conditions, we say that A is locally m-convex.

\begin{enumerate}
\def\labelenumi{\alph{enumi})}
\setcounter{enumi}{1}
\item
  Let A be a locally m-convex algebra. The map \((x,y)\mapsto xy\) The mapping from A × A into A is continuous. If A is unital, and if G denotes the set of invertible elements of A, then the map \(x\mapsto x^{-1}\) of G into A is continuous.
\item
  Let A be a unital locally m-convex algebra. The spectrum of every element of A is nonempty. If A is a field, then \(A = C \cdot 1\) .
\item
  Let A be a complete locally m-convex algebra. There exists a directed set I, a family \((\mathrm{A}_{i})_{i\in\mathrm{I}}\) of Banach algebras, continuous morphisms \(\pi_{ij}\colon A_{j}\mapsto A_{i}\) for \(i\leqslant j\) , with \(\pi_{ij}\circ\pi_{jk}=\pi_{ik}\) such that A is isomorphic to the subalgebra of \(\prod_{i}A_{i}\) consisting of \((x_{i})\) such that \(\pi_{ij}(x_{j})=x_{i}\) for \(i\leqslant j\) .
\item
  The algebra \(\mathcal{C}_{\mathbf{C}}(\mathbf{R})\) , equipped with the topology of compact convergence, is locally m-convex. The set of invertible elements of \(\mathcal{C}_{\mathbf{C}}(\mathbf{R})\) is dense in \(\mathcal{C}_{\mathbf{C}}(\mathbf{R})\) ; it is not open.
\item
  Let D be an open subset of C, and let A be the algebra of complex-valued holomorphic functions on D, equipped with the topology of compact convergence. Then A is locally m-convex. The set of invertible elements of A is closed.
\item
  The algebra of infinitely differentiable complex-valued functions on {[}0,1{]}, equipped with the topology of uniform convergence of each derivative, is locally m-convex.
\item
  Let A be the algebra of (classes of) complex functions f on {[}0,1{]} such that \(f \in \mathrm{L}^{p}([0,1])\) for every p \textgreater{} 1. Endowed with the topology defined by the family of seminorms \(f \mapsto \|f\|_{p}\) (p \textgreater{} 1), A is not locally m-convex. The map \((x,y) \mapsto xy\) from A × A into A is continuous.
\end{enumerate}

\exercisesection{Commutative Banach Algebras}

In the exercises below, all algebras considered are over C, unless explicitly stated otherwise.

\begin{enumerate}
\def\labelenumi{\arabic{enumi})}
\item
  Let A be a commutative Banach algebra and I a maximal ideal of A. Show that there are only two possible cases: either I is the kernel of a character of A, or I is a hyperplane containing \(A^{2}\) . (Use exerc. 6 of I, p.~154). In particular, if I is not closed, I is a dense hyperplane of A containing \(A^{2}\) . (To obtain an example of this latter situation, consider exerc. 3.)
\item
  Let A be the Banach algebra of sequences \(x = (x_{n})_{n \geqslant 1}\) of complex numbers tending to 0, with \(\|x\| = \sup_{i} |x_{i}|\) . Let I be the set of elements of A with finite support. Then I is a dense ideal of A, and is contained in no maximal ideal. (Observe that \(A^{2} = A\) , and use exerc. 1.)
\item
  Let \(\Lambda\) a set, \(p \in [1, +\infty[\) , and \(\mathbf{A} = \ell^p(\Lambda)\) the set of families \(x = (x_{\lambda})_{\lambda \in \Lambda}\) , where \(x_{\lambda} \in \mathbf{C}\) and
\end{enumerate}

\[
\| x \| _ {p} = \left(\sum_ {\lambda} | x _ {\lambda} | ^ {p}\right) ^ {1 / p} <   + \infty .
\]

\begin{enumerate}
\def\labelenumi{\alph{enumi})}
\item
  For componentwise addition and multiplication, and the norm \(x \mapsto \| x \|_p\) , the space A is a Banach algebra.
\item
  We have \(A^{2} \neq A\) and \(A^{2}\) is dense in A.
\item
  The space \(\mathsf{X}(\mathsf{A})\) of nonzero characters of A is naturally identified with \(\Lambda\) endowed with the discrete topology. (Use the fact that the dual of \(\ell^{p}(\Lambda)\) is \(\ell^{q}(\Lambda)\) , where q is the conjugate exponent of p.) The Gelfand transformation then becomes the identity map and its image in the space of continuous functions tending to 0 at infinity on \(\Lambda\) is dense and not closed.
\end{enumerate}

\begin{enumerate}
\def\labelenumi{\arabic{enumi})}
\setcounter{enumi}{3}
\tightlist
\item
  Let \((\alpha_{n})_{n\geqslant 0}\) a sequence in \(\mathbf{R}_{+}^{*}\) such that \(\alpha_0 = 1\) , \(\alpha_{m + n}\leqslant \alpha_m\alpha_n\) for all integers \(m\) and \(n\) , and such that \(\alpha_{n}^{1 / n}\to 0\) . Let A be the set of formal series \(x = \sum_{n = 0}^{\infty}x_{n}z^{n}\in \mathbf{C}[[z]]\) such that
\end{enumerate}

\[
\| x \| = \sum_ {n = 0} ^ {\infty} | x _ {n} | \alpha_ {n} <   + \infty .
\]

Show that A is a commutative unital Banach algebra topologically generated by z, and that the unique maximal ideal of A is the set of \(x \in A\) without constant term. (Observe that z is quasi-nilpotent.)

\begin{enumerate}
\def\labelenumi{\arabic{enumi})}
\setcounter{enumi}{4}
\item
  In the algebra \(\mathbf{C}(\mathrm{X})\) of rational fractions over \(\mathbf{C}\) , the ideal \(\{0\}\) is maximal, but is not of codimension 1.
\item
  Let \(\Delta\) the closed unit disk in C. Show that the group of invertible elements in \(\mathcal{C}(\Delta)\) is not dense in \(\mathcal{C}(\Delta)\) .
\item
  Let A be a commutative unital Banach algebra and \(\chi \in \mathsf{X}(\mathsf{A})\) . For every \(x \in \mathsf{A}\) , set \(\|x\|_{\chi} = |\chi(x)| + \|x - \chi(x)\|\) . Show that \(x \mapsto \|x\|_{\chi}\) is a norm, that \(\|xy\|_{\chi} \leqslant \|x\|_{\chi}\|y\|_{\chi}\) , \(\|e\|_{\chi} = 1\) , and that if \(\|e\| = 1\) , then \(\|x\| \leqslant \|x\|_{\chi} \leqslant 3\|x\|\) .
\item
  Let A be a commutative unital Banach algebra such that \(\|1\| = 1\) . Let N be the set of norms equivalent to the given norm, for which A is still a Banach algebra, and for which 1 has norm 1.
\end{enumerate}

\begin{enumerate}
\def\labelenumi{\alph{enumi})}
\tightlist
\item
  Let \(x_{0} \in A\) such that \(\varrho(x_{0}) < 1\) . For every \(x \in A\) , we set:
\end{enumerate}

\[
\left\| x \right\| ^ {\prime} = \inf \sum_ {n} \left\| a _ {n} \right\|
\]

for all representations of \(x\) in the form \(a_0 + a_1 x_0 + \cdots + a_n x_0^n\) (with \(a_0, a_1, \ldots, a_n \in \mathrm{A}\) ). Show that \(x \mapsto \|x\|^{\prime}\) is an element of N. (Since \(\varrho(x_0) < 1\) , there exists \(k\) such that \(\|x_0^n\| \leqslant k\) for all \(n\) , whence \(\|x\| \leqslant k \|x\|^{\prime}\) .)

\begin{enumerate}
\def\labelenumi{\alph{enumi})}
\setcounter{enumi}{1}
\item
  Show that \(\|x_{0}\|^{\prime}\leqslant1\) .
\item
  Deduce from a) and b) that, for every \(x \in A\) , one has
\end{enumerate}

\[
\varrho (x) = \inf _ {n \in \mathrm{N}} n (x).
\]

\begin{enumerate}
\def\labelenumi{\alph{enumi})}
\setcounter{enumi}{3}
\tightlist
\item
  If the only element of N is the given norm, we have \(A = C \cdot 1\) . (Use c) and exerc. 7.)
\end{enumerate}

\begin{enumerate}
\def\labelenumi{\arabic{enumi})}
\setcounter{enumi}{8}
\item
  Let X be a locally compact space and A the Banach algebra \(\mathcal{C}_{0}(X)\) . Let I be an ideal of A such that, for every \(x \in X\) , there exists \(f \in I\) such that \(f(x) \neq 0\) . Show that I contains the ideal of functions with compact support in X.
\item
  Let A be a commutative Banach algebra and D a continuous derivation of A into A, that is, a continuous linear map of A into A such that \(\mathrm{D}(ab)=(\mathrm{Da})b+a(\mathrm{Db})\) for \(a,b\in A\) .
\end{enumerate}

\begin{enumerate}
\def\labelenumi{\alph{enumi})}
\tightlist
\item
  For every \(\chi \in \mathsf{X}'(\mathbf{A})\) and every \(\lambda \in \mathbf{C}\) , the series
\end{enumerate}

\[
\varphi_ {\lambda} (a) = \sum_ {n = 0} ^ {\infty} \frac {\lambda^ {n}}{n !} \chi (\mathrm{D} ^ {n} (a))
\]

converges; \(\lambda \mapsto \varphi_{\lambda}(a)\) is an entire function, and \(a \mapsto \varphi_{\lambda}(a)\) is a character of A.

\begin{enumerate}
\def\labelenumi{\alph{enumi})}
\setcounter{enumi}{1}
\item
  The number \(\varphi_{\lambda}(a)\) is independent of \(\lambda\) . (Observe that \(|\varphi_{\lambda}(a)| \leqslant \|a\|\) according to \(a\) ).) Hence \(\chi(\mathrm{D}(a)) = 0\) .
\item
  The image of D is contained in the radical of A. Deduce from this a new proof of assertion f) of exerc. 31 of I, p.~162.
\end{enumerate}

\begin{enumerate}
\def\labelenumi{\arabic{enumi})}
\setcounter{enumi}{10}
\tightlist
\item
  Let B be the algebra over C of infinitely differentiable functions on {[}0,1{]} in C.
\end{enumerate}

\begin{enumerate}
\def\labelenumi{\alph{enumi})}
\tightlist
\item
  Let A be a subalgebra of B, equipped with a norm making A a Banach algebra. Show that there exists a sequence \((m_{n})_{n\geqslant0}\) into \(R_{+}\) such that, for all \(x\in A\) , one has
\end{enumerate}

\[
\sup _ {0 \leqslant t \leqslant 1} | x ^ {(n)} (t) | = \mathrm{O} (m _ {n}).
\]

(Let \(\mathrm{B}_n\) the Banach algebra of functions \(n\) times continuously differentiable on {[}0,1{]}. By prop. 9 of I, p.~40, the canonical injection \(\mathrm{A} \to \mathrm{B}_n\) is continuous; let \(\mathrm{M}_n\) be its norm; if \(x \in \mathrm{A}\) , one has \(\sup_{0 \leq t \leq 1} |x^{(n)}(t)| \leqslant n!\mathrm{M}_n \| x\|\) .)

\begin{enumerate}
\def\labelenumi{\alph{enumi})}
\setcounter{enumi}{1}
\tightlist
\item
  There exists no norm on B making it a Banach algebra.
\end{enumerate}

\begin{enumerate}
\def\labelenumi{\arabic{enumi})}
\setcounter{enumi}{11}
\tightlist
\item
  Let A be a unital Banach algebra and \((x_{n})\) a sequence of invertible elements of A tending to an element \(x \in A\) If the sequence \((\varrho(x_{n}^{-1}))\) is bounded, and if \(x_{n}x = xx_{n}\) for all n, then x is invertible. (By the cor. of prop. 8 of I, p.~38, \(\varrho(1 - x_{n}^{-1}x) \leqslant \varrho(x_{n}^{-1})\varrho(x_{n} - x)\) , hence \(x_{n}^{-1}x\) is invertible for n sufficiently large.)
\end{enumerate}

\begin{enumerate}
\def\labelenumi{¶ \arabic{enumi})}
\setcounter{enumi}{12}
\tightlist
\item
  Let A be the Banach algebra \(\ell^{2}(\mathbf{N})\) (exerc. 3). Let \(\mathrm{A}_{0}\) the subalgebra formed by sequences with finite support. Let B be the direct sum of \(\mathrm{A}_{0}\) and of \(\mathbf{C}\) , equipped with the multiplication \((f,\alpha)(g,\beta)=(fg,0)\) for \(f,g\in\mathrm{A}_{0},\alpha,\beta\in\mathbf{C}\) , and the norm
\end{enumerate}

\[
\left\| (f, \alpha) \right\| = \sup \left(\left\| f \right\|, \left| \alpha - \sum_ {n} f (n) \right|\right).
\]

\begin{enumerate}
\def\labelenumi{\alph{enumi})}
\item
  The completed algebra \(\widehat{\mathbf{B}}\) is a commutative Banach algebra whose radical R is \(\mathbf{C} \cdot (0,1)\) , and \(\widehat{\mathbf{B}}/\mathbf{R}\) is isometrically isomorphic to A.
\item
  Let \(\pi\colon\widehat{\mathrm{B}}\to\widehat{\mathrm{B}}/\mathrm{R}\) the canonical morphism. There exists no subalgebra \(\mathrm{B}_{1}\) of \(\widehat{\mathrm{B}}\) supplementary to R such that \(\pi|\mathrm{B}_{1}\colon\mathrm{B}_{1}\to\widehat{\mathrm{B}}/\mathrm{R}=\mathrm{A}\) is a homeomorphism. (Let \(\mathrm{B}_{1}\) be such a subalgebra. Let \(u_{n}\in\mathrm{A}\) such that \(u_{n}(n)=1\) , \(u_{n}(m)=0\) for \(m\neq n\) . Let \(e_{n}\in\mathrm{B}_{1}\) such that \(\pi(e_{n})=u_{n}\) . Show that \(e_{n}\in\mathrm{A}_{0}\) , that \(\sum_{n=1}^{\infty}\frac{1}{n}u_{n}\) converges in A but \(\sum_{n=1}^{\infty}\frac{1}{n}e_{n}\) does not converge in \(\widehat{\mathrm{B}}\) .)
\end{enumerate}

\begin{enumerate}
\def\labelenumi{\arabic{enumi})}
\setcounter{enumi}{13}
\tightlist
\item
  In \(\mathbf{C}^2\) , let \(\Omega\) the compact set defined by
\end{enumerate}

\[
\frac {1}{2} \leqslant | z _ {1} | ^ {2} + | z _ {2} | ^ {2} \leqslant 1,
\]

and let A be the subalgebra of \(\mathcal{C}(\Omega)\) consisting of functions that are holomorphic in \(\mathring{\Omega}\) . Let U be the open set defined by \(|z_1|^2 + |z_2|^2 < 1\) . By virtue of a theorem of Hartogs (cf.~J.-P. DEMAILLY, Complex analytic and differential geometry, Ch. I, th. 3.28), for every \(f \in \mathrm{A}\) , there exists a function \(\widetilde{f} \in \mathcal{C}(\overline{\mathrm{U}})\)

, unique, which is holomorphic in U and coincides with f on \(\Omega\) . Show that the canonical injection from A into \(\mathcal{C}(\Omega)\) is an isometric morphism h whose image is a full subalgebra of \(\mathcal{C}(\Omega)\) but \(\mathsf{X}(h)\) is not surjective.

\begin{enumerate}
\def\labelenumi{\arabic{enumi})}
\setcounter{enumi}{14}
\item
  Let A and B be commutative unital Banach algebras, \(\varphi\) a unital morphism from A to B and \((x_{\lambda})_{\lambda \in \Lambda}\) a family of elements of A such that the closed full subalgebra of A generated by the \(x_{\lambda}\) is equal to A. For \(\mathsf{X}(\varphi)\) to be surjective, it is necessary and sufficient that \(\mathrm{Sp}_{\mathrm{A}}^{\Lambda}((x_{\lambda})) = \mathrm{Sp}_{\mathrm{B}}^{\Lambda}((\varphi (x_{\lambda})))\) (Use diagram (1) of no. 6 of I, p.~41, where the right-hand arrow is bijective by prop. 12 of I, p.~43.)
\item
  Let A be a commutative Banach algebra. The following conditions are equivalent:
\end{enumerate}

\begin{enumerate}
\def\labelenumi{(\roman{enumi})}
\item
  A is radical-free and \(\mathcal{G}(\mathrm{A})\) is closed in \(\mathcal{C}_0(\mathsf{X}(\mathrm{A}))\) ;
\item
  \(\mathcal{G}\) is a homeomorphism from A onto \(\mathcal{G}(\mathrm{A})\) ;
\item
  there exists a constant c \textgreater{} 0 such that \(\|x\|^{2} \leqslant c\|x^{2}\|\) for all \(x \in A\) .
\end{enumerate}

\begin{enumerate}
\def\labelenumi{\arabic{enumi})}
\setcounter{enumi}{16}
\item
  Let A be a commutative Banach algebra and I a closed ideal of A. Denote by h the canonical injection of I into A. The space \(\mathsf{X}'(\mathsf{I})\) is identified with the quotient space of \(\mathsf{X}'(\mathsf{A})\) by the equivalence relation defined by \(\mathsf{X}'(h)\) . (Use part (a) of Exercise 4 of I, p.~153).
\item
  Let A be a unital commutative Banach algebra and \(x\) an element of A. Suppose that the closed full subalgebra of A generated by \(x\) is equal to A, so that the map \(\chi \mapsto \chi(x)\) allows one to identify \(\mathsf{X}(\mathsf{A})\) and \(\mathrm{Sp}_{\mathsf{A}}(x)\) . Let \(\mathcal{H}\) the set of functions \(\mathcal{G}y\) on \(\mathrm{Sp}_{\mathsf{A}}(x)\) , for \(y\) ranging over A. Then \(\check{\mathrm{S}}_{\mathcal{H}}(\mathrm{Sp}_{\mathsf{A}}(x))\) (INT, IV, §7, no. 4) is the boundary of \(\mathrm{Sp}_{\mathsf{A}}(x)\) relative to C. (To show that \(\check{\mathrm{S}}_{\mathcal{H}}(\mathrm{Sp}_{\mathsf{A}}(x))\) is contained in this boundary, use the maximum principle. Conversely, let \(z_0\) a point of this boundary; to show that \(z_0 \in \check{\mathrm{S}}_{\mathcal{H}}(\mathrm{Sp}_{\mathsf{A}}(x))\) , consider \((x - z_1)^{-1}\) for \(z_1\) quite close to \(z_0\) into \(\mathbf{C} = \mathrm{Sp}_{\mathsf{A}}(x)\) .)
\item
  Let A be a commutative unital Banach algebra, B a closed unital subalgebra of A, and T the map \(\chi \mapsto \chi|B\) from X(A) into X(B), and let R be the equivalence relation on X(A) defined by T. Then T induces, by passing to the quotient, a homeomorphism from X(A)/R onto T(X(A)), and for every \(x \in B\) , the function \(f = \mathcal{G}_{\mathrm{B}}(x)\) is such that \(|f|\) attains its upper bound on T(X(A)).
\item
  Let \(\Delta\) the disk \(|z| \leqslant 1\) into \(\mathbf{C}\) the set of \(f \in \mathcal{C}(\Delta)\) that are holomorphic in the interior of \(\Delta\) and \(\mathrm{A}_1\) the Banach subalgebra of \(\mathcal{C}(\Delta)\) generated by A and by the function \(z \mapsto |z|\) . Show that \(\mathrm{X}(\mathrm{A}_1)\) is homeomorphic to the set of \((x_{1}, x_{2}, x_{3}) \in \mathbf{R}^{3}\) such that \(x_{1}^{2} + x_{2}^{2} \leqslant x_{3}^{2}\) and \(0 \leqslant x_{3} \leqslant 1\) .
\item
  \begin{enumerate}
  \def\labelenumii{\alph{enumii})}
  \tightlist
  \item
    Let \((\mathrm{X}_{\lambda})_{\lambda\in\Lambda}\) a family of indeterminates. For any formal series \(f\in\mathbf{C}[[(\mathrm{X}_{\lambda})]]_{\lambda\in\Lambda}\) , let \(\|f\|\) the sum of the absolute values of its coefficients. Let A be the set of \(f\in\mathbf{C}[[(\mathrm{X}_{\lambda})]]_{\lambda\in\Lambda}\) such that \(\|f\|<+\infty\) Equipped with the usual addition and multiplication, A is a commutative unital Banach algebra without radical.
  \end{enumerate}
\end{enumerate}

\begin{enumerate}
\def\labelenumi{\alph{enumi})}
\setcounter{enumi}{1}
\tightlist
\item
  For every commutative unital Banach algebra B, there exists an algebra A of the type in a) and a continuous unital morphism from A onto B.
\end{enumerate}

\begin{enumerate}
\def\labelenumi{\arabic{enumi})}
\setcounter{enumi}{21}
\tightlist
\item
  Let \(\Omega\) a compact subset of \(\mathbf{C}^n\) . The map
\end{enumerate}

\[
(z _ {1}, \ldots , z _ {n}) \longmapsto (z _ {1}, \ldots , z _ {n}, \overline {{z}} _ {1}, \ldots , \overline {{z}} _ {n})
\]

is a homeomorphism from \(\Omega\) onto a compact polynomially convex subset of \(\mathbf{C}^{2n}\) . (Let \(A = \mathcal{C}(\Omega)\) , \(z_i\) the coordinate functions on \(\mathbf{C}^n\) . Then \((z_i|\Omega, \overline{z}_i|\Omega)_{1 \leqslant i \leqslant n}\) is a topological generating system of A, and the map in question is the map from \(X(A) = \Omega\) into \(\mathbf{C}^{2n}\) defined by this topological generating system.)

\begin{enumerate}
\def\labelenumi{\arabic{enumi})}
\setcounter{enumi}{22}
\tightlist
\item
  We denote \(\Omega\) the set of \((z_1, z_2) \in \mathbf{C}^2\) such that \(|z_1| \leqslant 1\) , \(|z_2| \leqslant 1\) . Let \(\alpha \in [0,1]\) and \(\Omega_\alpha\) the set of \((z_1, z_2) \in \mathbf{C}^2\) such that
\end{enumerate}

\[
| z _ {1} | \leqslant 1, \quad | z _ {2} | \leqslant (1 - \alpha) | z _ {1} | + \alpha .
\]

\begin{enumerate}
\def\labelenumi{\alph{enumi})}
\item
  The complements of \(\Omega\) and \(\Omega_{\alpha}\) into \(\mathbf{C}^2\) are connected.
\item
  \(\Omega\) is the polynomially convex hull of \(\Omega_{\alpha}\) . (If \((z_1^0, z_2^0) \in \Omega\) and if \(\mathrm{P} \in \mathbf{C}[z_1, z_2]\) , there exists \(z \in \mathbf{C}\) such that \(|z| = 1\) and \(|\mathrm{P}(z_1^0, z_2^0)| \leqslant |\mathrm{P}(z, z_2^0)|\) ; but \(|\mathrm{P}(z, z_2^0)| \leqslant \sup_{(z_1, z_2) \in \Omega_{\alpha}} |\mathrm{P}(z_1, z_2)|\) .)
\end{enumerate}

\begin{enumerate}
\def\labelenumi{\arabic{enumi})}
\setcounter{enumi}{23}
\item
  Let \(n \geqslant 1\) . There exists a closed convex subset of \(\mathbf{C}^{n}\) that is not polynomially convex.
\item
  Let A be a complete commutative unital locally m-convex algebra (cf.~I, p.~165, exerc. 37).
\end{enumerate}

\begin{enumerate}
\def\labelenumi{\alph{enumi})}
\item
  An element x of A is invertible if and only if \(\chi(x) \neq 0\) for every continuous character \(\chi\) of A. The radical of A is the intersection of the kernels of the continuous characters of A, hence is closed.
\item
  In example e) of exercise 37 of I, p.~165, let I be the ideal formed by the f such that \(f(n) = 0\) for every integer n \textgreater{} 0 sufficiently large. Then every maximal ideal containing I is dense and of infinite codimension.
\end{enumerate}

\begin{enumerate}
\def\labelenumi{\arabic{enumi})}
\setcounter{enumi}{25}
\tightlist
\item
  Let X be a compact topological space and E a complex Banach space, whose norm is denoted by \(\|\cdot\|_{E}\) . We denote by \(\mathcal{C}(X,E)\) the complex vector space of continuous functions from X into E.
\end{enumerate}

\begin{enumerate}
\def\labelenumi{\alph{enumi})}
\tightlist
\item
  Equipped with the norm
\end{enumerate}

\[
\| f \| = \sup _ {x \in \mathrm{X}} \| f (x) \| _ {\mathrm{E}},
\]

the space \(\mathcal{C}(\mathrm{X},\mathrm{E})\) is a Banach space.

\begin{enumerate}
\def\labelenumi{\alph{enumi})}
\setcounter{enumi}{1}
\tightlist
\item
  Let B be the vector subspace of \(\mathcal{C}(\mathrm{X},\mathrm{E})\) generated by the functions \(x\mapsto f(x)y\) , where \(f\in \mathcal{C}(\mathrm{X})\) and \(y\in \mathrm{E}\) . Show that B is dense in \(\mathcal{C}(\mathrm{X},\mathrm{E})\) .
\end{enumerate}

In what follows, we consider a commutative Banach algebra A.

\begin{enumerate}
\def\labelenumi{\alph{enumi})}
\setcounter{enumi}{2}
\tightlist
\item
  Equipped with the multiplication \((fg)(x) = f(x)g(x)\) for all \(x\in \mathrm{X}\) , the Banach space \(\mathcal{C}(\mathrm{X},\mathrm{A})\) is a commutative Banach algebra.
\end{enumerate}

\begin{enumerate}
\def\labelenumi{\alph{enumi})}
\setcounter{enumi}{3}
\tightlist
\item
  Every character \(\chi\in\mathsf{X}(\mathcal{C}(\mathrm{X},\mathrm{A}))\) is of the form \(f\mapsto\widetilde{\chi}(f(x))\) , where \(\widetilde{\chi}\in\mathsf{X}(\mathrm{A})\) and \(x\in\mathrm{X}\) . (First treat the case where A is unital and \(\|1\|=1\) ; consider the restriction of \(\chi\) to \(\mathcal{C}(\mathrm{X})\cdot1\) and to the image of the homomorphism \(\mathrm{A}\to\mathcal{C}(\mathrm{X},\mathrm{A})\) which sends \(a\) to the constant function \(x\mapsto a\) ).
\end{enumerate}

\begin{enumerate}
\def\labelenumi{\alph{enumi})}
\setcounter{enumi}{4}
\tightlist
\item
  The map from \(\mathsf{X}(\mathrm{A})\times \mathrm{X}\) into \(\mathsf{X}(\mathcal{C}(\mathrm{X},\mathrm{A}))\) which to \((\widetilde{\chi},x)\) associates the character \(f\mapsto \widetilde{\chi} (f(x))\) is a homeomorphism.
\end{enumerate}

\begin{enumerate}
\def\labelenumi{\arabic{enumi})}
\setcounter{enumi}{26}
\tightlist
\item
  Let A be a commutative unital Banach algebra over C. A closed subset F of X(A) is called a boundary of A if one has
\end{enumerate}

\[
\sup _ {\chi \in \mathsf {X} (\mathrm{A})} | \chi (x) | = \sup _ {\chi \in \mathrm{F}} | \chi (x) |
\]

for all \(x \in A\) .

\begin{enumerate}
\def\labelenumi{\alph{enumi})}
\item
  Order the set of boundaries of A by inclusion. There exists a minimal boundary.
\item
  Let S be a minimal boundary of A. Every boundary of A contains S, hence S is the intersection of all boundaries of A. (Let F be a boundary of A; show that for every \(x \in S\) , every neighborhood U of x meets F; for this, use the fact that S --- U is not a boundary of A.)
\end{enumerate}

One says that S is the Shilov boundary of A, and one denotes it by \(\partial A\) .

\begin{enumerate}
\def\labelenumi{\alph{enumi})}
\setcounter{enumi}{2}
\item
  If the image of the Gelfand transform of A is dense in \(\mathcal{C}(\mathrm{X}(\mathrm{A}))\) , the Shilov boundary of A is equal to \(\mathrm{X}(\mathrm{A})\) .
\item
  Let A be the algebra of continuous functions on the unit disk \(\Delta\) of \(z \in C\) such that \(|z| \leqslant 1\) that are holomorphic inside of \(\Delta\) (example 9 of I, p.~20). Then \(\mathsf{X}(\mathsf{A})\) is identified with \(\Delta\) (exercise 6 of I, p.~193) and the Shilov boundary of A is identified with the unit circle U.
\end{enumerate}

\begin{enumerate}
\def\labelenumi{\arabic{enumi})}
\setcounter{enumi}{27}
\item
  Let X be a compact topological space and n an integer \(\geqslant 1\) . Let I be a two-sided ideal of the Banach algebra \(\mathcal{C}(\mathrm{X},\mathrm{M}_{n}(\mathbf{C}))\) . Show that there exists a unique closed subset S ⊂ X such that I is the set of \(f \in \mathcal{C}(\mathrm{X}, \mathrm{M}_{n}(\mathbf{C}))\) whose restriction to S is zero.
\item
  * Let X be a compact real manifold. We equip \(\mathcal{C}^{\infty}(\mathrm{X})\) of the topology of uniform convergence of functions and all their derivatives; it is a Fréchet space. Let A be a subalgebra of \(\mathcal{C}(\mathrm{X})\) We assume that A is equipped with a norm \(f \mapsto \| f\|_{\mathrm{A}}\) which makes it a Banach algebra, such that the canonical injection of A into \(\mathcal{C}(\mathrm{X})\) is continuous. We further assume that \(\mathcal{C}^{\infty}(\mathrm{X})\) is contained and dense in A.
\end{enumerate}

\begin{enumerate}
\def\labelenumi{\alph{enumi})}
\tightlist
\item
  There exists a constant \(C \geqslant 0\) and an integer \(k \in N\) such that for all \(f \in \mathcal{C}^{\infty}(X)\) , one has
\end{enumerate}

\[
\| f \| _ {\mathrm{A}} \leqslant \operatorname * {C} \sup _ {x \in \mathrm{X}} | f ^ {(k)} (x) |.
\]

\begin{enumerate}
\def\labelenumi{\alph{enumi})}
\setcounter{enumi}{1}
\item
  The subalgebra A is full in \(\mathcal{C}(\mathrm{X})\) . (Let \(f \in \mathrm{A}\) that does not vanish on X and \(\delta = \inf_{x \in \mathrm{X}} |f(x)|\) ; choose \(g \in \mathcal{C}^{\infty}(\mathrm{X})\) such that \(\|f - g\|_{\mathrm{A}} < \delta/3\) , then note that \(f^{-1} = g^{-1} \sum_{n \in \mathbf{N}} ((g - f)/g)^{n}\) into \(\mathcal{C}(\mathrm{X})\) and verify that the series converges in A.)
\item
  Deduce from this a new proof of Wiener's theorem (I, p.~38, example).
\end{enumerate}

\exercisesection{Holomorphic Functional Calculus}

In the exercises below, all algebras considered are over C, unless explicitly stated otherwise.

\begin{enumerate}
\def\labelenumi{\arabic{enumi})}
\tightlist
\item
  Let \(n \geqslant 1\) be an integer and \(A = \mathcal{L}(\mathbf{C}^{n})\) . Let \(\alpha \in C\) and \(x \in A\) an element whose matrix relative to the canonical basis of \(C^{n}\) is a Jordan matrix of order n and eigenvalue \(\alpha\) The spectrum of x is reduced to \(\alpha\) .
\end{enumerate}

\begin{enumerate}
\def\labelenumi{\alph{enumi})}
\tightlist
\item
  For \(f \in \mathcal{O}(\{\alpha\})\) compute the matrix relative to the canonical basis of \(\mathbf{C}^n\) which represents \(f(x)\) .
\end{enumerate}

Let \(m \geqslant 1\) be an integer. Let \(B_{m}\) the unital Banach algebra of jets of functions of class \(C^{m}\) in a neighborhood of \(\alpha\) endowed with the topology of \(C^{m}\) -uniform convergence (VAR, R2, §12).

\begin{enumerate}
\def\labelenumi{\alph{enumi})}
\setcounter{enumi}{1}
\item
  There exists a continuous morphism of unital algebras from \(B_{m-1}\) to A such that \(\varphi(z)=x\) , where \(z\in B_{m-1}\) is the jet of order m-1 at \(\alpha\) of the identity function of C.
\item
  There does not exist a continuous morphism \(\varphi\) of unital algebras from \(B_m\) to A such that \(\varphi(z) = x\) , where \(z \in B_m\) is the jet of the identity function of \(\mathbf{C}\) .
\end{enumerate}

\begin{enumerate}
\def\labelenumi{\arabic{enumi})}
\setcounter{enumi}{1}
\tightlist
\item
  Let A be a unital Banach algebra and G the group of invertible elements of A.
\end{enumerate}

\begin{enumerate}
\def\labelenumi{\alph{enumi})}
\item
  For an element \(x \in A\) to be of the form \(\exp(y)\) , where \(y \in A\) , it is necessary and sufficient that it be contained in a connected commutative subgroup of G.
\item
  For an element \(x \in \mathbf{A}\) to be of the form \(\exp(y)\) , where \(y \in \mathbf{A}\) , it suffices that 0 belong to the unbounded connected component of \(\mathbf{C} = \mathrm{Sp}_{\mathrm{A}}(x)\) .
\end{enumerate}

\begin{enumerate}
\def\labelenumi{\arabic{enumi})}
\setcounter{enumi}{2}
\tightlist
\item
  Let A be a unital Banach algebra, G the group of invertible elements of A, and \(G_{1}\) the identity component of G.
\end{enumerate}

\begin{enumerate}
\def\labelenumi{\alph{enumi})}
\item
  Let \(x \in \mathbf{A}\) and \(n\) an integer \(\geqslant 0\) such that \(x^n = e\) . Show that \(x \in \mathbf{G}_1\) . (The spectrum of \(x\) is finite. The set of \(\lambda \in \mathbf{C}\) such that \(y(\lambda) = \lambda x + e - \lambda e\) such that it is invertible has finite complement, hence is connected. Now, \(y(0) = e\) , \(y(1) = x\) .)
\item
  Assume A is commutative. Every element of G/G₁ distinct from the identity element has infinite order. (If \(x \in G\) and \(x^n \in G_1\) , one has \(x^n = \exp(y)\) with a \(y \in A\) by no. 10 of I, p.~78, whence ( \(x \exp(-y/n)^n = e\) . Apply a).)
\end{enumerate}

\begin{enumerate}
\def\labelenumi{\arabic{enumi})}
\setcounter{enumi}{3}
\item
  There exists a commutative Banach algebra, non-unital and non-zero, such that \(X(A)\) is compact.
\item
  Let A be a unital Banach algebra and \(a \in A\) . Let G be the group of invertible elements of A and \(G_1\) the connected component of the identity in G. The exponential spectrum of A is the set of \(\lambda \in \mathbf{C}\) such that \(x - \lambda e\) does not belong to \(G_1\) (cf.~prop. 12 of I, p.~79).
\end{enumerate}

\begin{enumerate}
\def\labelenumi{\alph{enumi})}
\item
  The exponential spectrum of a is closed; it contains the spectrum of a.
\item
  The boundary of the exponential spectrum of a is contained in the singular spectrum of a (exercise 16 of I, p.~158).
\item
  The exponential spectrum of a is the union of the spectrum of a and of certain bounded connected components of \(\mathbf{C}=\mathrm{Sp}(a)\) .
\end{enumerate}

\begin{enumerate}
\def\labelenumi{\alph{enumi})}
\setcounter{enumi}{3}
\tightlist
\item
  Let \(f\in \mathbf{C}[\mathrm{X}]\) be a polynomial. The image under \(f\) of the exponential spectrum of \(a\) contains the exponential spectrum of \(f(a)\) . In general, the inclusion is strict (use Fredholm theory, cf.~III, to appear).
\end{enumerate}

\begin{enumerate}
\def\labelenumi{\arabic{enumi})}
\setcounter{enumi}{5}
\tightlist
\item
  Let A be a commutative unital Banach algebra.
\end{enumerate}

\begin{enumerate}
\def\labelenumi{\alph{enumi})}
\item
  If j is an idempotent of A, one has \(\exp(2i\pi j)=1\) and \(\exp(i\pi j)=1-2j\) .
\item
  Let \(p \in A\) such that \(\exp(p) = 1\) . Then \(\mathrm{Sp}_{\mathrm{A}}(p)\) consists of a finite number of points of the form \(2in\pi\) , where \(n \in Z\) .
\item
  For \(\lambda \notin \mathrm{Sp}_{\mathrm{A}}(p)\) , we have:
\end{enumerate}

\[
\int_ {0} ^ {1} \exp (t (p - \lambda 1)) d t = (1 - \exp (- \lambda)) (\lambda 1 - p) ^ {- 1}.
\]

\begin{enumerate}
\def\labelenumi{\alph{enumi})}
\setcounter{enumi}{3}
\tightlist
\item
  Deduce that every point of \(\mathrm{Sp}_{\mathbf{A}}(p)\) is a simple pole of the resolvent \(\mathrm{R}(p,\lambda)\) , then that there exists an integer \(m \geqslant 0\) , integers \(n_{k} \in Z\) and pairwise orthogonal idempotents \(j_{k}\) , for \(1 \leqslant k \leqslant m\) , such that
\end{enumerate}

\[
p = 2 i \pi \sum_ {k} n _ {k} j _ {k}.
\]

\begin{enumerate}
\def\labelenumi{¶ \arabic{enumi})}
\setcounter{enumi}{6}
\tightlist
\item
  a) Let B be a complex Banach space and j a point of B such that \(\|j\| = 1\) . For every \(x \in B\) , the limit
\end{enumerate}

\[
\lim_{\substack{\alpha \to 0\\ \alpha >0}}\frac{1}{\alpha} (\| j + \alpha x\| -1)
\]

exists. Let \(\varphi(x)\) be this limit.

\begin{enumerate}
\def\labelenumi{\alph{enumi})}
\setcounter{enumi}{1}
\tightlist
\item
  Let
\end{enumerate}

\[
\psi (x) = \sup_{\substack{\varepsilon \in \mathbf{C}\\ |\varepsilon | = 1}}\varphi (\varepsilon x).
\]

The function \(\psi\) is a semi-norm bounded above by the norm of B. One has

\[
\psi (x) = \sup _ {f} | f (x) |,
\]

where \(f\) ranges over the set of elements of the dual of B such that \(\|f\| = f(j) = 1\) .

\begin{enumerate}
\def\labelenumi{\alph{enumi})}
\setcounter{enumi}{2}
\tightlist
\item
  One now assumes that B is a Banach algebra and that \(j = 1\) is a unit element of B. One has:
\end{enumerate}

\[
\varphi (x) = \lim_{\substack{\alpha \to 0\\ \alpha >0}}\frac{1}{\alpha}\log (\| \exp (\alpha x)\|) = \sup_{\alpha >0}\frac{1}{\alpha}\log (\| \exp (\alpha x)\|).
\]

(Observe that \(\log(\|\exp((\alpha+\alpha')x)\|) \leqslant \log(\|\exp(\alpha x)\|) + \log(\|\exp(\alpha' x)\|).\) )

\begin{enumerate}
\def\labelenumi{\alph{enumi})}
\setcounter{enumi}{3}
\tightlist
\item
  Using c), show that \(\|\exp(\lambda x)\| \leqslant \exp(|\lambda|\psi(x))\) . By integrating \(\lambda \mapsto \lambda^{-2} \exp(\lambda x)\) over the circle determined by \(|\lambda| = \varrho\) , deduce that
\end{enumerate}

\[
\| x \| \leqslant \varrho^ {- 1} \exp (\varrho \psi (x)),
\]

whence, setting \(\varrho = \psi(x)^{-1}\) , that \(\|x\| \leqslant e\psi(x)\) , where \(e = \exp(1)\) .

\begin{enumerate}
\def\labelenumi{\alph{enumi})}
\setcounter{enumi}{4}
\tightlist
\item
  For every \(x\) nonzero element of B, there exists a linear form \(f \in \mathrm{B}'\) such that \(f(1) = \| f \| = 1\) and \(f(x) \neq 0\) .
\end{enumerate}

\begin{enumerate}
\def\labelenumi{\arabic{enumi})}
\setcounter{enumi}{7}
\item
  Let A be a unital Banach algebra, \(x \in A\) , and B the subalgebra of A generated by x. For \(\mathrm{Sp}_{\mathrm{A}}(x)\) to be a finite set all of whose elements are poles of the resolvent of x, it is necessary and sufficient that B be of finite dimension. (To show that the condition is necessary, if \(\mathrm{R}(x, \lambda) = \sum_{i=1}^{p} \mathrm{F}_{i}(\lambda) a_{i}\) , where \(a_{i} \in A\) and where \(F_{i}\) is a scalar rational function, the expansion of \(\mathrm{R}(x, \lambda)\) around \(\lambda = \infty\) shows that the \(x^{n}\) are linear combinations of the \(a_{i}\) .)
\item
  Let A be a unital Banach algebra, \(x \in A\) , U an open neighborhood of \(\mathrm{Sp}_{\mathrm{A}}(x)\) having only a finite number of connected components \(U_{i}\) ( \(i \in I\) ), and \(f \in \mathcal{O}(\mathrm{U})\) . For \(f(x) = 0\) , it is necessary and sufficient that the following conditions be satisfied:
\end{enumerate}

\begin{enumerate}
\def\labelenumi{(\roman{enumi})}
\item
  for every i such that \(\mathrm{Sp}_{\mathrm{A}}(x)\cap\mathrm{U}_{i}\) is infinite or contains a point which is not a pole of the resolvent, one has \(f|U_{i}=0\) ;
\item
  every pole of order \(p\) of the resolvent is a zero of order \(\geqslant p\) of \(f\) (For necessity of (ii), argue as in Exerc. 8.)
\end{enumerate}

\begin{enumerate}
\def\labelenumi{\arabic{enumi})}
\setcounter{enumi}{9}
\tightlist
\item
  Let \(\Delta\) the disk \(|z| \leqslant 1\) in C. Let A be the space of functions f that are continuous in \(\Delta\) and holomorphic in the interior of \(\Delta\) and such that there exists a sequence \((c_{n})\) summable of complex numbers satisfying
\end{enumerate}

\[
f (z) = \sum_ {n \geqslant 0} c _ {n} z ^ {n}
\]

for all \(z \in \Delta\) . We then denote by \(\|f\| = \sum_{n}|c_{n}|\) .

\begin{enumerate}
\def\labelenumi{\alph{enumi})}
\item
  Equipped with the multiplication \((f,g)\mapsto fg\) , the space A is a unital Banach algebra and z is a topological generator of A.
\item
  The space \(\mathsf{X}(\mathsf{A})\) is identified with \(\Delta\) ; if \(f \in A\) and if g is holomorphic in an open neighborhood of \(f(\Delta)\) , then \(g \circ f \in A\) .
\end{enumerate}

\begin{enumerate}
\def\labelenumi{\arabic{enumi})}
\setcounter{enumi}{10}
\item
  Let A be a commutative unital Banach algebra and \(x \in A\) . Let S be a subset of X(A) closed for the Jacobson topology, and let f be a complex holomorphic function in a neighborhood of \((\mathcal{G}x)(\mathrm{S})\) . There exists \(y \in A\) such that \(Gy = f \circ Gx\) on S. (Let I be the intersection of the kernels of the \(\chi \in S\) . Use the functional calculus in A/I.)
\item
  Let U be a nonempty open set in C. Show that there exists no norm on \(\mathcal{O}(\mathrm{U})\) for which \(\mathcal{O}(\mathrm{U})\) is a Banach algebra. (Using th. 1 of I, p.~29, show that the functions of \(\mathcal{O}(\mathrm{U})\) would be bounded.)
\item
  Let \(n\in \mathbf{N}\)
\end{enumerate}

\begin{enumerate}
\def\labelenumi{\alph{enumi})}
\tightlist
\item
  A subset V of \(C^{n}\) is polynomially convex if and only if, for every \(x \in C^{n} - V\) , there exists an entire function \(f: C^{n} \to C\) such that
\end{enumerate}

\[
| f (x) | > \sup _ {y \in V} | f (y) |.
\]

\begin{enumerate}
\def\labelenumi{\alph{enumi})}
\setcounter{enumi}{1}
\item
  Let \(K \subset R^{n}\) a compact set. Then K is polynomially convex in \(C^{n}\) .
\item
  Let V be a polynomially convex compact subset of \(C^{n}\) and \(f \in \mathcal{O}(V)\) . Then the graph of f is polynomially convex in \(C^{n+1}\) . (Use the Oka--Weil theorem.)
\end{enumerate}

\begin{enumerate}
\def\labelenumi{\arabic{enumi})}
\setcounter{enumi}{13}
\item
  Let K be a compact subset of C and let P (resp. Q) be the closure in \(\mathcal{O}(\mathrm{K})\) of the set of germs of polynomial (resp. rational) functions holomorphic in a neighborhood of K. Show that \(\mathrm{P} = \mathrm{Q}\) if and only if the set \(\mathbf{C} - \mathbf{K}\) is connected.
\item
  Let U be the set of complex numbers of modulus 1. Let \(K_{1}\) , \(K_{2}\) be the compact subsets of \(C^{2}\) given by
\end{enumerate}

\[
\mathrm{K} _ {1} = \mathbf {U} \times \{0 \}, \quad \mathrm{K} _ {2} = \{(z, \overline {{z}}) \mid z \in \mathbf {U} \}.
\]

\begin{enumerate}
\def\labelenumi{\alph{enumi})}
\item
  Show that the set of germs in a neighborhood of \(K_{1}\) of polynomial functions on C with complex coefficients is not dense in \(\mathcal{O}(K_{1})\) .
\item
  Show that the set of germs in a neighborhood of \(K_{2}\) of polynomial functions on C with complex coefficients is dense in \(\mathcal{O}(K_{2})\) .
\end{enumerate}

\begin{enumerate}
\def\labelenumi{\arabic{enumi})}
\setcounter{enumi}{15}
\tightlist
\item
  Let A be a unital Banach algebra, \(x \in A\) , and k an integer \(\geqslant 0\) . Let U be the set of complex numbers of modulus 1.
\end{enumerate}

\begin{enumerate}
\def\labelenumi{\alph{enumi})}
\tightlist
\item
  Suppose that \(\mathrm{Sp}_{\mathrm{A}}(x)\subset\mathbf{U}\) and that \(\|\mathrm{R}(x,\lambda)\|=\mathrm{O}((1-|\lambda|)^{k})\) when \(|\lambda|\) tends to 1. Show that \(\|x^{n}\|=\mathrm{O}(|n|^{k})\) when n tends to \(\pm\infty\) . (For every \(\delta>1\) , show that
\end{enumerate}

\[
2 \pi \| x ^ {n} \| \leqslant \mathrm{M} \frac {\delta^ {n}}{(\delta - 1) ^ {k}} + \mathrm{M} \frac {\delta^ {- n}}{(1 - \delta^ {- 1}) ^ {k}}
\]

where M is independent of n and \(\delta\) . Take \(\delta = n/(n - k)\) when n tends to \(+\infty\) , \(\delta = n/(n + k)\) when n tends to \(-\infty\) .

\begin{enumerate}
\def\labelenumi{\alph{enumi})}
\setcounter{enumi}{1}
\tightlist
\item
  Suppose that \(\|x^{n}\|=\mathrm{O}(|n|^{k})\) when n tends to \(\pm\infty\) . Show that \(\operatorname{Sp}_{\mathbf{A}}(x)\subset\mathbf{U}\) and that \(\|\operatorname{R}(\lambda,x)\|=\operatorname{O}((1-|\lambda|)^{k+1})\) when \(|\lambda|\) tends to 1. (We have \(\varrho(x)=\varrho(x^{-1})=1\) , hence \(\operatorname{Sp}_{\mathbf{A}}(x)\subset\mathbf{U}\) . For \(|\lambda|<1\) ,
\end{enumerate}

\[
\mathrm{R} (\lambda , x) = - x ^ {- 1} \left(1 + \lambda x ^ {- 1} + \lambda^ {2} x ^ {- 2} + \lambda^ {3} x ^ {- 3} + \dots\right).
\]

Deduce an upper bound for \(\|\mathrm{R}(\lambda,x)\|\) , and argue similarly if \(|\lambda|>1\) .

\begin{enumerate}
\def\labelenumi{\arabic{enumi})}
\setcounter{enumi}{16}
\tightlist
\item
  Let \(X_{1}\) (resp. \(X_{2}\) ) the set of \((z_{1}, z_{2}) \in \mathbf{C}^{2}\) such that
\end{enumerate}

\[
\frac {1}{2} \leqslant | z _ {1} | ^ {2} + | z _ {2} | ^ {2} \leqslant 1, \quad \text { resp. } \quad | z _ {1} | ^ {2} + | z _ {2} | ^ {2} \leqslant 1.
\]

Let \(X = X_{1} \cup \{(0,0)\} \subset \mathbf{C}^{2}\) . Let \(A = \mathcal{C}(X)\) , and \(a_{1}, a_{2}\) the restrictions to X of the coordinate functions on \(C^{2}\) .

\begin{enumerate}
\def\labelenumi{\alph{enumi})}
\item
  The simultaneous spectrum \(\mathrm{Sp}_{\mathrm{A}}^{2}(a_{1},a_{2})\) is equal to X.
\item
  Construct two distinct continuous unital morphisms from \(\mathcal{O}(\mathrm{X})\) to A sending \(z_{1}\) to \(a_{1}\) and \(z_{2}\) to \(a_{2}\) . (Use the following theorem: for every function \(f\) holomorphic in a neighborhood of \(\mathrm{X}_{1}\) , there exists a function g holomorphic in a neighborhood of \(X_{2}\) which coincides with f in a neighborhood of \(X_{1}\) (cf. I, p. 168, exerc. 14).
\end{enumerate}

\begin{enumerate}
\def\labelenumi{¶ \arabic{enumi})}
\setcounter{enumi}{17}
\tightlist
\item
  Let A be a commutative unital Banach algebra.
\end{enumerate}

\begin{enumerate}
\def\labelenumi{\alph{enumi})}
\tightlist
\item
  For any finite family I of elements of A, let S(I) ⊂ C \(^{I}\) the simultaneous spectrum of I. If I′ extends I, the canonical surjection pr \(_{I,I'}\) of C \(^{I'}\) on C \(^{I}\) is such that pr \(_{I,I'}\) (S(I')) = S(I), whence an injective morphism \(\varphi_{II'}\) of \(\mathcal{O}(S(I))\) into \(\mathcal{O}(S(I'))\) . Let \(\mathcal{O}(X(A))\) the inductive limit of the \(\mathcal{O}(S(I))\) with respect to the \(\varphi_{II'}\) (we index the finite families of elements of A by the finite subsets of A × N). The canonical maps
\end{enumerate}

\[
\varphi_ {\mathrm{I}} \colon \mathcal {O} (\mathrm{S} (\mathrm{I})) \longrightarrow \mathcal {O} (\mathrm{X} (\mathrm{A}))
\]

are injective.

We shall equip \(\mathcal{O}(\mathrm{X(A)})\) of the inductive limit topology of those of \(\mathcal{O}(\mathrm{S(I)})\) . For every \(a \in \mathrm{A}\) the coordinate function in a neighborhood of \(\mathrm{S}(a) \subset \mathbf{C}\) defines an element \(z_a\) of \(\mathcal{O}(\mathrm{X(A)})\) .

\begin{enumerate}
\def\labelenumi{\alph{enumi})}
\setcounter{enumi}{1}
\tightlist
\item
  The continuous surjection \(\mathsf{X}(\mathsf{A}) \to \mathsf{S}(\mathsf{I})\) defines a continuous morphism \(\mathcal{C}(\mathsf{S}(\mathsf{I})) \to \mathcal{C}(\mathsf{X}(\mathsf{A}))\) . The morphisms \(\mathcal{O}(\mathsf{S}(\mathsf{I})) \to \mathcal{C}(\mathsf{S}(\mathsf{I})) \to \mathcal{C}(\mathsf{X}(\mathsf{A}))\) define a continuous morphism:
\end{enumerate}

\[
\varphi \colon \mathcal {O} (\mathrm{X} (\mathrm{A})) \longrightarrow \mathcal {C} (\mathrm{X} (\mathrm{A})).
\]

This morphism is not injective in general.

\begin{enumerate}
\def\labelenumi{\alph{enumi})}
\setcounter{enumi}{2}
\item
  There exists a unique continuous unital morphism \(\psi\) of \(\mathcal{O}(\mathsf{X}(\mathsf{A}))\) to A such that \(\psi(z_a) = a\) for all \(a \in \mathsf{A}\) (Use th. 1 of I, p.~51.) Its composite with the Gelfand transformation is \(\varphi\) .
\item
  Let E be a separated locally convex space, U a weakly open subset of E, and \(f: U \to C\) a function. We say that f is weakly holomorphic if, for every \(u \in U\) there exists a weakly open neighborhood \(V_u\) of u, a closed vector subspace \(F_u\) of finite codimension in E, and a holomorphic function \(g_u\) on \(p_u(V_u)\) , where \(p_u\) denotes the canonical mapping of E onto \(E/F_u\) , such that \(f|V_u = g_u \circ p_u\) .
\end{enumerate}

Show that if \(f\) is weakly holomorphic, and if K is a weakly compact subset of U, there exists a weakly open neighborhood \(V_K\) of K in U, a closed vector subspace \(F_K\) of finite codimension in E, and a holomorphic function \(g_K\) on \(p_K(V_K)\) such that \(f|V_K = g_K \circ p_K\) , where \(p_K\) is the canonical mapping of E onto \(E/F_K\) .

\begin{enumerate}
\def\labelenumi{\alph{enumi})}
\setcounter{enumi}{4}
\tightlist
\item
  Show that \(\mathcal{O}(\mathsf{X}(\mathsf{A}))\) is the inductive limit of the \(\mathcal{O}(\mathsf{S})\) , where S ranges over the set of weakly open neighborhoods of \(\mathsf{X}(\mathsf{A})\) in the dual of A, and where \(\mathcal{O}(\mathsf{S})\) denotes the algebra of weakly holomorphic functions in S, endowed with the topology of compact convergence.
\end{enumerate}

\exercisesection{Regular Commutative Banach Algebras}

In the exercises below, all algebras considered are over C.

\begin{enumerate}
\def\labelenumi{\arabic{enumi})}
\item
  Let A be a regular commutative Banach algebra and I a closed ideal of A. Then I and A/I are regular commutative Banach algebras.
\item
  Let A be a regular commutative Banach algebra without radical, I an ideal of A, \(\chi\) an element of V(I), and J the set of \(x \in A\) such that Gx has compact support disjoint from \(\{\chi\}\) . Then I + J is the smallest of the ideals R of A such that R ⊃ I and V(R) = \(\{\chi\}\) .
\item
  Let A be a commutative Banach algebra, A' the dual Banach space of A. For every \(x \in A\) , let \(\kappa(x)\) the linear operator in A' transposed to multiplication by x in A. For \(x \in A\) and \(x' \in A'\) , we denote by \(x \star x' = \kappa(x)(x') \in A'\) . A vector subspace of A' is said to be invariant if it is stable under all the \(\kappa(x)\) .
\end{enumerate}

\begin{enumerate}
\def\labelenumi{\alph{enumi})}
\item
  For an element \(x'\) of A' to be proportional to a nonzero character of A, it is necessary and sufficient that \(Cx'\) be invariant.
\item
  The map \(V \mapsto V^{\circ}\) is a bijection from the set of closed ideals of A onto the set of weakly closed invariant vector subspaces of \(A'\) .
\item
  If W is an invariant vector subspace of A', we denote by \(\sigma(\mathrm{W})\) the set of characters belonging to W. If \(W_{1}\) and \(W_{2}\) are weakly closed invariant subspaces, and if W is the weakly closed vector subspace generated by \(W_{1}\) and \(W_{2}\) , one has \(\sigma(\mathrm{W}) = \sigma(\mathrm{W}_{1}) \cup \sigma(\mathrm{W}_{2})\) (Use b).
\item
  Let W be a weakly closed invariant vector subspace of A′. Then \(\sigma(\mathrm{W})\) is the set of \(\chi \in \widehat{A}\) such that the condition \(x * x' = 0\) for all \(x' \in W\) implies \(\chi(x) = 0\) .
\item
  If A admits a unit element, every weakly closed invariant vector subspace of \(A'\) contains a nonzero character.
\item
  In the remainder of this exercise, we assume that A is a regular algebra without radical and with identity. If W is a weakly closed invariant vector subspace of A' and U a neighborhood of \(\sigma(\mathrm{W})\) into \(\mathsf{X}(\mathsf{A})\) then W is contained in the weakly closed vector subspace of A' generated by U.
\end{enumerate}

\begin{enumerate}
\def\labelenumi{\alph{enumi})}
\setcounter{enumi}{6}
\tightlist
\item
  Let \(x' \in A'\) and let W be the weakly closed invariant vector subspace of \(A'\) generated by \(x'\) . We set \(\sigma(x') = \sigma(\mathrm{W})\) . If \(x \in A\) is such that \(\mathcal{G}x = 1\) on a neighborhood of \(\sigma(x')\) , then \(x \star x' = x'\) . (Use \(f\) ).)
\end{enumerate}

\begin{enumerate}
\def\labelenumi{\alph{enumi})}
\setcounter{enumi}{7}
\tightlist
\item
  If \(\sigma(x')\) is the union of two disjoint closed sets \(\sigma_1\) and \(\sigma_2\) , there exists \(x_1', x_2' \in A'\) such that \(\sigma(x_1') = \sigma_1\) , \(\sigma(x_2') = \sigma_2\) and \(x' = x_1' + x_2'\) . (Set \(x_1' = u_1 \star x'\) , \(x_{2}^{\prime}=u_{2}\star x^{\prime}\) , with \(\mathcal{G}u_{1}=1\) (resp. 0) in a neighborhood of \(\sigma_{1}\) (resp. \(\sigma_{2}\) ), and \(\mathcal{G}u_{2}=1\) (resp. 0) in a neighborhood of \(\sigma_{2}\) (resp. \(\sigma_{1}\) ).)
\end{enumerate}

\begin{enumerate}
\def\labelenumi{\arabic{enumi})}
\setcounter{enumi}{3}
\item
  Let A be a commutative regular semisimple Banach algebra satisfying Ditkin's condition, and let I be a closed ideal of A. \(x \in \Upsilon(\mathrm{V}(\mathrm{I}))\) and F the boundary of \(\mathrm{V}(Ax)\) . If \(\mathrm{F} \cap \mathrm{V}(\mathrm{I})\) contains no nonempty perfect set, we have \(x \in I\) (Imitate the reasoning of Prop. 5 of I, p.~94.)
\item
  Let A be a regular commutative Banach algebra without radical. We assume that every closed ideal of A is an intersection of regular maximal ideals. Let \(x \in A\) . Let F be the set of zeros of \(\mathcal{G}'(x)\) Then x is a limit of elements \(u_{n}x\) , where \(u_{n} \in A\) and \(\mathcal{G}(u_{n})\) vanishes in a neighborhood of F. (Let J be the set of \(y \in A\) such that \(\mathcal{G}(y)\) with compact support disjoint from F. We have \(x \in \overline{J}\) . Use prop. 4 of I, p.~91.) In particular, A satisfies Ditkin's condition.
\item
  Let B be the algebra of continuously differentiable complex functions on {[}0, 1{]}, equipped with the norm \(\|f\| = \sup |f| + \sup |f'|\) . Let A be the closed subalgebra of B formed by the functions \(f \in \mathrm{B}\) such that \(f\left(\frac{1}{2}\right) = 0\) . Let I be the closed ideal of A formed by the \(f \in \mathrm{A}\) such that \(f'\left(\frac{1}{2}\right) = 0\) .
\end{enumerate}

\begin{enumerate}
\def\labelenumi{\alph{enumi})}
\item
  The space X(A) is identified with {[}0, 1{]} = \{ \(\frac{1}{2}\) \}, and A is regular.
\item
  The ideal I is a maximal ideal of A that is not regular.
\end{enumerate}

\begin{enumerate}
\def\labelenumi{\arabic{enumi})}
\setcounter{enumi}{6}
\tightlist
\item
  Let R, \(T_{1}\) , \(T_{2}\) be the subsets formed by \(C^{4}\) consisting of the elements \((z_{1}, z_{2}, z_{3}, z_{4}) \in \mathbf{C}^{4}\) such that:
\end{enumerate}

\[
\mathrm{R} \colon z _ {1} z _ {2} = 2, \qquad 1 \leqslant | z _ {1} | \leqslant 2, \qquad z _ {3} = z _ {4} = 0,
\]

\[
\mathrm{T} _ {1} \colon z _ {1} z _ {2} = 2, \qquad | z _ {1} | = 1, \qquad | z _ {3} | \leqslant 1, \qquad z _ {4} = 0,
\]

\[
\mathrm{T} _ {2} \colon z _ {1} z _ {2} = 2, \qquad | z _ {1} | = 2, \qquad | z _ {3} | \leqslant 1, \qquad z _ {4} = z _ {3} ^ {2}.
\]

\begin{enumerate}
\def\labelenumi{\alph{enumi})}
\tightlist
\item
  The set \(X = R \cup T_{1} \cup T_{2}\) is polynomially convex. (Show that X is the set of \((z_{1}, z_{2}, z_{3}, z_{4})\) such that
\end{enumerate}

\[
z _ {1} z _ {2} = 2, \quad | z _ {1} | \leqslant 2, \quad | z _ {2} | \leqslant 2, \quad | z _ {3} | \leqslant 1, \quad z _ {4} (z _ {4} - z _ {3} ^ {2}) = 0,
\]

\[
| z _ {4} z _ {2} ^ {k} | \leqslant 1 \quad \text { and } \quad | z _ {1} ^ {k} (z _ {4} - z _ {3} ^ {2}) | \leqslant 1 \quad \text { for } \quad k = 1, 2, \ldots).
\]

\begin{enumerate}
\def\labelenumi{\alph{enumi})}
\setcounter{enumi}{1}
\item
  Let A be the closed subalgebra of \(\mathcal{C}(\mathrm{X})\) generated by the restrictions to X of polynomial functions on \(\mathbf{C}^4\) . Then X is identified with \(\mathsf{X}(\mathsf{A})\) .
\item
  Let \(\Gamma_1, \Gamma_2 \subset \mathbf{C}\) the circles centered at 0 with radii 1, 2 oriented in the positive direction. Let \(\mu_0, \mu_1, \nu\) the measures on \(\Gamma_1, \Gamma_2, \Gamma_1\) defined by the differential forms \(z^{-1} dz, -z^{-1} dz, z^{-2} dz\) . Let \(\mu = \mu_0 + \mu_1\) , so that \(\mu \otimes \nu\) is a measure on \((\Gamma_1 \cup \Gamma_2) \times \Gamma_1\) . Let B be the set of \(z = (z_1, z_2, z_3, z_4) \in X\) such that \((z_1, z_3) \in (\Gamma_1 \cup \Gamma_2) \times \Gamma_1\) . The map \(\varphi: z \mapsto (z_1, z_3)\) from B into \((\Gamma_{1} \cup \Gamma_{2}) \times \Gamma_{1}\) is a homeomorphism. Let \(\nu\) the measure \(\varphi^{-1}(\mu \otimes \nu)\) . Show that \(\nu\) is orthogonal to A.
\item
  Let g be the function on X that is zero on \(R \cup T_{1}\) and equal to \(z_{3}\) on \(T_{2}\) . Then \(g \notin A\) .
\item
  The function g coincides with an element of A in a neighborhood of every \(z \in X\) .
\end{enumerate}

\exercisesection{Stellar Algebras}

\begin{enumerate}
\def\labelenumi{\arabic{enumi})}
\item
  Find a unital involutive Banach algebra A and a unitary element \(u \in A\) such that \(\|u\| \neq 1\) .
\item
  Let A be a Banach algebra and \(A_{1} \subset A\) a closed subalgebra. If \(A_{1}\) is an involutive algebra, then every character of \(A_{1}\) extends to a character of A.
\item
  Let A be a unital involutive algebra. Suppose that for every x in A, the spectrum of x does not contain -1.
\end{enumerate}

\begin{enumerate}
\def\labelenumi{\alph{enumi})}
\item
  For every idempotent element x of A, show that there exists a hermitian idempotent element y of A such that \(xA = yA\) . (Consider the inverse z of \(1 + (x - x^{*})(x^{*} - x)\) ; show that z commutes with x and set \(y = xx^{*}z\) .)
\item
  For every idempotent element \(x\) of A, show that there exists a hermitian idempotent element \(y\) of A and \(u\) invertible in A such that \(x = uyu^{-1}\) . (If \(x\mathrm{A} = y\mathrm{A}\) , set \(u = (1 - (y - x))^{-1}\) .)
\end{enumerate}

\begin{enumerate}
\def\labelenumi{\arabic{enumi})}
\setcounter{enumi}{3}
\item
  Let A be a normed algebra, equipped with a continuous involution. If one sets \(\|x\|' = \sup(\|x\|, \|x^{*}\|)\) for all \(x \in A\) , A becomes an involutive normed algebra, and the new norm is equivalent to the old one.
\item
  Let A be a commutative Banach algebra without radical. Every involution of A is continuous. (Use prop. 9 of I, p.~40.)
\end{enumerate}

\begin{enumerate}
\def\labelenumi{¶ \arabic{enumi})}
\setcounter{enumi}{5}
\tightlist
\item
  Let A be a commutative unital involutive Banach algebra whose every character is hermitian. Suppose that, for every \(x \in A\) non-invertible hermitian element, one has \(\|(x - \lambda)^{-1}\| = o(\mathcal{I}(\lambda)^{-2})\) when \(\lambda\) tends to 0 with \(\mathcal{I}(\lambda) \neq 0\) . Then every closed ideal I of A is an intersection of maximal ideals. (Let \(\varphi: A \longrightarrow A/I\) be the canonical morphism. Let \(x \in A\) such that \(\varphi(x)\) be in the radical of A/I. It must be proved that \(\varphi(x) = 0\) . Reduce to the case where \(x\) is hermitian. Then use exerc. 29 of I, p.~161).
\end{enumerate}

\begin{enumerate}
\def\labelenumi{\arabic{enumi})}
\setcounter{enumi}{6}
\item
  Let A be a commutative unital involutive Banach algebra. For every character of A to be hermitian, it is necessary and sufficient that \(1 + xx^{*}\) be invertible for every \(x \in A\) . (If every character of A is hermitian, \(\mathcal{G}(1 + xx^{*})\) is everywhere \textgreater{} 0 on \(\mathsf{X}(\mathsf{A})\) . If a character \(\chi\) of A is non-hermitian, there exists a \(x \in A\) hermitian element such that \(\chi(x) = i\) , whence \(\chi(1 + xx^{*}) = 0\) and \(1 + xx^{*}\) is not invertible.)
\item
  Construct an example of a group G such that the involutive Banach algebra \(\mathcal{M}^{1}(G)\) (I, p.~100, example 4) is not a star algebra.
\item
  Let A be a unital Banach algebra. The group of unitary elements of A is a maximal bounded subgroup of A*.
\item
  Let A be an involutive Banach algebra and \(\widetilde{\mathsf{A}}\) the algebra obtained by adjoining a unit element. Show that the norm defined in number I, p.~15 on \(\widetilde{\mathsf{A}}\) does not coincide, in general, with the norm defined by prop. 3 of I, p.~105.
\item
  Let A be an involutive Banach algebra, whose involution satisfies \(\|x^{*}x\| = \|x\| \cdot \|x^{*}\|\) for all \(x \in A\) . Then A is a star algebra. (One has \(\|y^{2}\| = \|y\|^{2}\) for hermitian y, whence \(\|y\| = \varrho(y)\) ; observing that \(\mathrm{Sp}_{\mathrm{A}}^{\prime}(x^{*}) = \overline{\mathrm{Sp}_{\mathrm{A}}^{\prime}(x)}\) for all \(x \in A\) , we therefore have
\end{enumerate}

\[
\| x ^ {*} \| \cdot \| x \| = \| x ^ {*} x \| = \varrho (x ^ {*} x) \leqslant \varrho (x ^ {*}) \varrho (x) = \varrho (x ^ {2}) \leqslant \| x \| ^ {2},
\]

whence \(\|x^{*}\|\leqslant\|x\|\) and \(\|x^{*}\|=\|x\|\) .

\begin{enumerate}
\def\labelenumi{\arabic{enumi})}
\setcounter{enumi}{11}
\item
  Let A be a star algebra, N the set of normal elements of A, \(x \in N\) and V a neighborhood of 0 in C. There exists a neighborhood U of x in N such that, for every \(y \in U\) , one has \(\mathrm{Sp}_{\mathrm{A}}(y) \subset \mathrm{Sp}_{\mathrm{A}}(x) + \mathrm{V}\) and \(\mathrm{Sp}_{\mathrm{A}}(x) \subset \mathrm{Sp}_{\mathrm{A}}(y) + \mathrm{V}\) . (Use the cor. of 1 of I, p.~108, prop. 3 of I, p.~23 and prop. 10 of I, p.~76.)
\item
  Let A be a unital Banach algebra over C, \(x \in A\) , and \(S = \mathrm{Sp}_{\mathrm{A}}(x)\) . Suppose that there exists a morphism \(\varphi\) of \(\mathcal{C}(\mathrm{S})\) into A which sends 1 to 1 and z to x, denoting by z the identity map of S. Then \(\|\varphi(f)\| \geqslant \|f\|\) for every \(f \in \mathcal{C}(\mathrm{S})\) . In particular, if the morphism \(\varphi\) is continuous, then \(\varphi\) is a homeomorphism onto its image. (One may assume A commutative. Let \(\lambda \in S\) . There exists \(\chi \in \mathrm{X}(\mathrm{A})\) such that \(\lambda = \chi(x) = \chi(\varphi(z)) = (\mathrm{X}(\varphi)(\chi))(z)\) . Therefore \(\mathrm{X}(\varphi)(\chi)\) be the character \(f \mapsto f(\lambda)\) of \(\mathcal{C}(\mathrm{S})\) . Let \(f \in \mathcal{C}(\mathrm{S})\) . According to what precedes, there exists \(\chi \in \mathrm{X}(\mathrm{A})\) such that \(|f|\) attains its maximum at \(\mathrm{X}(\varphi)(\chi)\) . Then \(|\chi(\varphi(f))| = \|f\|\) , hence \(\|\varphi(f)\| \geqslant \|f\|\) .)
\item
  Let G be a locally compact group. If Stell(G) has an identity element, then the group G is discrete. (Use the fact that, in the Hilbert space \(L^{2}(G)\) , the identity map is the norm limit of endomorphisms \(\gamma_{f}\) , where \(\gamma\) is the left regular representation and where \(f \in L^{1}(G)\) .)
\item
  \begin{enumerate}
  \def\labelenumii{\alph{enumii})}
  \tightlist
  \item
    Let A be a commutative stellar algebra. If x and y are positive elements of A such that \(x \leqslant y\) , then \(x^{2} \leqslant y^{2}\) .
  \end{enumerate}
\end{enumerate}

\begin{enumerate}
\def\labelenumi{\alph{enumi})}
\setcounter{enumi}{1}
\item
  There exists a stellar algebra A and elements x and y of A such that \(0 \leqslant x \leqslant y\) but such that \(x^{2}\) is not \(\leqslant y^{2}\) .
\item
  There exists a stellar algebra A and elements \(x\) and \(y\) in A such that \(|x + y|\) is not \(\leqslant |x| + |y|\) .
\end{enumerate}

\begin{enumerate}
\def\labelenumi{\arabic{enumi})}
\setcounter{enumi}{15}
\tightlist
\item
  Let X be a topological space. Let \(\mathcal{C}_{b}(\mathrm{X})\) be the commutative stellar algebra of bounded continuous complex-valued functions on X.
\end{enumerate}

\begin{enumerate}
\def\labelenumi{\alph{enumi})}
\item
  Let \(\widehat{\mathrm{X}} = \mathrm{X}(\mathcal{C}_b(\mathrm{X}))\) ; the space \(\widehat{\mathrm{X}}\) is a compact space such that \(\mathcal{C}(\widehat{\mathrm{X}})\) is identified with \(\mathcal{C}_b(\mathrm{X})\) .
\item
  The map \(j: X \to \widehat{X}\) which sends x to the character \(f \mapsto f(x)\) is continuous.
\item
  If X is completely regular (TG, IX, p.~8, def. 4), the map j is a homeomorphism of X onto a dense subspace of \(\widehat{X}\) (To show that the image of j is dense in \(\widehat{X}\) , show that a continuous complex function on \(\widehat{X}\) vanishing on \(j(\mathrm{X})\) is identically zero.)
\item
  The space \(\widehat{X}\) is identified with the Stone–Čech compactification of X (TG, IX, p.~9).
\end{enumerate}

\begin{enumerate}
\def\labelenumi{\arabic{enumi})}
\setcounter{enumi}{16}
\item
  Let A be a C*-algebra and \(x \in \mathrm{A}\) . If \(x = a - b\) with \(a, b\) positive and \(ab = ba = 0\) , then \(a = x^{+}\) and \(b = x^{-}\) .
\item
  Let A be a separable stellar algebra. There exists a sequence \((e_{n})_{n\in\mathbf{N}}\) into \(A_{+}\) such that the associated elementary filter (TG, I, p.~42, def. 7) is an increasing approximate unit of A.
\item
  Let A be a stella algebra. An element \(x\) is said to be systematically positive if \(x \in \mathrm{A}_{+}\) and if \(x\mathrm{A}x\) is dense in A.
\end{enumerate}

\begin{enumerate}
\def\labelenumi{\alph{enumi})}
\item
  Show that a commutative stellar algebra A contains a systematically positive element if, and only if, there exists a locally compact topological space X, countable at infinity, such that A is isomorphic to \(\mathcal{C}_{0}(X)\) .
\item
  If A is unital, then an element x of A is strictly positive if and only if x is positive and invertible.
\item
  Let \(x \in A\) an element that is systematically positive such that \(\|x\| \leqslant 1\) Prove that the sequence \((x^{1/n})_{n \geqslant 1}\) is an increasing approximate unit in A.
\item
  A stellate algebra A contains a systematically positive element if and only if there exists a sequence \((x_{n})_{n\geqslant1}\) in A, which forms an approximate identity for A.
\end{enumerate}

\begin{enumerate}
\def\labelenumi{\arabic{enumi})}
\setcounter{enumi}{19}
\item
  Let A be a C*-algebra, I a closed two-sided ideal of A, and B a closed involutive subalgebra of A. Then B+I is a closed involutive subalgebra of A, and there exists an isometric isomorphism \(B + I/I \rightarrow B/B \cap I\) .
\item
  Let S be the closed unit disk in C. Let A be the involutive Banach algebra of continuous functions from S to C that are holomorphic in the open unit disk, where the involution is given by \(f^{*}(z) = \overline{f(\overline{z})}\) There exists a Hermitian element of A whose spectrum is not contained in R.
\item
  Let \(\mathrm{A} = \mathcal{C}([0,1])\) Show that the ideal I of A generated by the identity function of {[}0, 1{]} is not closed.
\item
  Let A be a unital C*-algebra. We denote by \(A_{+}^{\prime}\) the vector space of linear forms f on A, not necessarily continuous, such that
\end{enumerate}

\[
f (\mathrm{A} _ {+}) \subset [ 0, + \infty [
\]

(“positive linear forms”).

\begin{enumerate}
\def\labelenumi{\alph{enumi})}
\item
  Let \(f \in A_{+}^{\prime}\) . For every sequence \((x_{n})_{n \in \mathbf{N}}\) into \(A_{+}^{\leqslant 1}\) , the family \((f(x_{n}))\) is bounded. (Show that the series \(\sum_{i} f(x_{i})y_{i}\) converges for every positive family \((y_{i})\) into \(\ell^{1}(\mathbf{N})\) .)
\item
  The linear map f is continuous.
\item
  A linear form \(f\) belongs to \(\mathrm{A}_{+}^{\prime}\) if and only if \(f\) is continuous and \(\| f \| = f(1)\) .
\item
  Let \(x \in A_{h}\) a Hermitian element. There exists a positive linear form \(f \in A_{+}^{\prime}\) such that \(|f(x)| = \|x\|\) (Use Exercise 7 of I, p.~174).
\end{enumerate}

\begin{enumerate}
\def\labelenumi{\arabic{enumi})}
\setcounter{enumi}{23}
\tightlist
\item
  Let A be a unital Banach algebra over C. Suppose that \(x \mapsto x^{*}\) and \(x \mapsto x^{\circ}\) are involutions on A such that A, equipped with each one, is a stellar algebra. We denote by \(A_{1}\) and \(A_{2}\) these two stellar algebras.
\end{enumerate}

\begin{enumerate}
\def\labelenumi{\alph{enumi})}
\item
  Using the notation from Exercise 23, show that \(\mathrm{A}_{1,+}^{\prime} = \mathrm{A}_{2,+}^{\prime}\) . One denotes by \(\mathrm{A}_{+}^{\prime}\) this set.
\item
  Let \(x \in A_{1,h}\) an element such that \(x^{*} = x\) . Show that for every \(f \in A_{+}^{\prime}\) , one has \(f(i(x - x^{\circ})) = 0\) . Deduce that \(x \in A_{2,h}\) (Use part (b) of Exercise 23).
\item
  Conclude that \(A_{1} = A_{2}\) , and that a complex Banach algebra admits at most one involution making it a C*-algebra.
\end{enumerate}

\begin{enumerate}
\def\labelenumi{\arabic{enumi})}
\setcounter{enumi}{24}
\tightlist
\item
  Let G be a group. We denote by C{[}G{]} the group algebra of G over C. It is a subspace of the Hilbert space \(\ell^{2}(\mathrm{G})\) . For g in G, we denote by {[}g{]} the element of C{[}G{]} corresponding to g. The elements {[}g{]}, for \(g \in G\) , form an orthonormal basis of \(\ell^{2}(G)\) and C{[}G{]} is dense in \(\ell^{2}(G)\) .
\end{enumerate}

\begin{enumerate}
\def\labelenumi{\alph{enumi})}
\tightlist
\item
  The map
\end{enumerate}

\[
\sum_ {g} \alpha_ {g} [ g ] \mapsto \sum_ {g} \overline {{\alpha}} _ {g} [ g ^ {- 1} ]
\]

is an involution on \(\mathbf{C}[\mathrm{G}]\) .

\begin{enumerate}
\def\labelenumi{\alph{enumi})}
\setcounter{enumi}{1}
\tightlist
\item
  For \(g \in G\) , let \(\pi(g)\) the map \(x \mapsto gx\) of C{[}G{]}. The map which associates \(\pi(g)\) to g extends to a continuous and injective morphism of involutive algebras from C{[}G{]} into \(\mathcal{L}(\ell^{2}(G))\) .
\end{enumerate}

We denote \(\operatorname{Stell}_{r}(\mathbf{G})\) , and the reduced C*-algebra of G is called the closure in \(\mathcal{L}(\ell^{2}(\mathbf{G}))\) of the image of the map \(\pi\) . It is a C*-algebra. We identify \(\mathbf{C}[\mathbf{G}]\) with its image under \(\pi\) into \(\operatorname{Stell}_{r}(\mathbf{G})\) .

\begin{enumerate}
\def\labelenumi{\alph{enumi})}
\setcounter{enumi}{2}
\tightlist
\item
  The linear trace map \(\operatorname{Tr} \colon \mathbf{C}[\mathrm{G}] \to \mathbf{C}\) such that \(\operatorname{Tr}(1) = 1\) and \(\operatorname{Tr}(g) = 0\) for \(g \neq 1\) extends to a continuous linear map \(\operatorname{Stell}_r(\mathrm{G}) \to \mathbf{C}\) .
\end{enumerate}

\begin{enumerate}
\def\labelenumi{\alph{enumi})}
\setcounter{enumi}{3}
\tightlist
\item
  We have \(\operatorname{Tr}(xy) = \operatorname{Tr}(yx)\) for all \(x\) and \(y\) into \(\operatorname{Stell}_r(G) \) .
\end{enumerate}

\begin{enumerate}
\def\labelenumi{\alph{enumi})}
\setcounter{enumi}{4}
\item
  We have \(\operatorname{Tr}(xx^{*}) \geqslant 0\) for all \(x\) into \(\operatorname{Stell}_r(G)\) , and \(\operatorname{Tr}(xx^{*}) = 0\) if and only if \(x = 0\) .
\item
  Let e be an idempotent element of C{[}G{]} such that \(e \neq 0\) and \(e \neq 1\) . Show that \(0 < \operatorname{Tr}(e) < 1\) . (Use Exercise 3).
\item
  Let e be an idempotent element of Z{[}G{]}. Show that e = 0 or e = 1. ("Kaplansky's theorem").
\item
  Let x be an element of C{[}G{]}. Then x is invertible if and only if it is left-invertible, if and only if it is right-invertible.
\item
  Let n be a strictly positive integer. A matrix x of type \((n, n)\) with coefficients in C{[}G{]} is invertible if and only if it is left-invertible, if and only if it is right-invertible.
\end{enumerate}

\begin{enumerate}
\def\labelenumi{\arabic{enumi})}
\setcounter{enumi}{25}
\item
  Let A be a unital C*-algebra. Let \(a\) and \(b\) be normal elements of A and \(x \in \mathrm{A}^*\) such that \(b = xax^{-1}\) . There exists a unitary element \(u\) of A such that \(b = uau^*\) . (Use the preceding exercise and the polar decomposition of \(x\) , noting that \(|x|\) belongs to the bicommutant of \(x\) ).
\item
  Let A be a unital C*-algebra, and let a, b, and c be elements of A. Suppose that b and c are normal. If ac = cb, then we have \(f(a)c = cf(b)\) for every continuous function f on the union of the spectrum of a and the spectrum of b.
\item
  Let A be a C*-algebra. Let \(a\) and \(b\) be positive elements of A such that \(ab \neq 0\) . Then \(\mathrm{Sp}_{\mathrm{A}}(ab)\) is not reduced to 0.
\item
  \begin{enumerate}
  \def\labelenumii{\alph{enumii})}
  \tightlist
  \item
    Let G be a locally compact topological group. The Banach algebra L \(^{1}\) (G) has an approximate identity. (Use INT, VIII, § 2, no. 7, lemma 4).
  \end{enumerate}
\end{enumerate}

\begin{enumerate}
\def\labelenumi{\alph{enumi})}
\setcounter{enumi}{1}
\item
  Let A be a Banach algebra admitting an approximate identity. For every character \(\chi\in\mathsf{X}(\mathrm{A})\) , one has \(\|\chi\|=1\) .
\item
  The Banach subalgebra of L \(^{1}\) (Z) formed by the sequences ( \(x_{n}\) such that \(x_{n} = 0\) if \(n < 0\) has no approximate identity. (Construct a nonzero character of norm \(< 1\) .)
\end{enumerate}

\begin{enumerate}
\def\labelenumi{\alph{enumi})}
\setcounter{enumi}{3}
\tightlist
\item
  Let G be the free group generated by two elements \(a\) and \(b\) and \(\mu\) a Haar measure on G. Let I be the closed ideal of \(\mathrm{L}^1 (\mathrm{G})\) consisting of the \(f\in \mathrm{L}^{1}(\mathrm{G})\) such that \(\mu (f) = 0\) . Let \(\xi = e^{2i\pi /3}\) . Let \(x = \varepsilon_{e} + \xi \varepsilon_{a} + \xi^{-1}\varepsilon_{b}\in \mathrm{L}^{1}(\mathrm{G})\) . For every \(f\in \mathrm{L}^{1}(\mathrm{G})\) , one has \(\| x*f - x\| \geqslant 1\) The Banach algebra I has no approximate identity.
\end{enumerate}

\begin{enumerate}
\def\labelenumi{¶ \arabic{enumi})}
\setcounter{enumi}{29}
\tightlist
\item
  Let A be a Banach algebra with an approximate identity. Denote by \(\widetilde{\mathsf{A}}\) the Banach algebra obtained from A by adjoining an identity element, denoted by 1.
\end{enumerate}

\begin{enumerate}
\def\labelenumi{\alph{enumi})}
\item
  Let \(x \in \mathrm{A}\) . Then \(x\) belongs to the closed left ideal generated by \(x\) .
\item
  For every \(\varepsilon > 0\) and every finite family \((x_i)_{i \in \mathrm{I}}\) of elements of \(\mathrm{A}\) , there exists \(e \in \mathrm{A}\) such that \(\| e \| \leqslant 1\) and \(\| ex_i - x_i \| < \varepsilon\) for all \(i \in \mathrm{I}\) .
\item
  Let \(x \in A\) , \(\varepsilon > 0\) and \(0 < \gamma \leqslant 1/8\) . There exists \(\eta > 0\) such that for all \(e \in A\) satisfying \(\|e\| \leqslant 1\) and \(\|ex - x\| < \eta\) , then \(\alpha = (1 - \gamma) \cdot 1 + \gamma e\) is invertible in \(\widetilde{A}\) and \(\|x - \alpha^{-1}x\| < \varepsilon\) .
\end{enumerate}

\begin{enumerate}
\def\labelenumi{\alph{enumi})}
\setcounter{enumi}{3}
\tightlist
\item
  Let \(x \in \mathrm{A}\) , \(\varepsilon > 0\) , \(\delta > 0\) and \(0 < \gamma \leqslant 1/8\) . There exists a sequence \((e_n)_{n \geqslant 1}\) in A having the following properties:

\begin{enumerate}
\def\labelenumi{(\roman{enumi})}
\item
  We have \(\|e_{n}\| < 2\) for all \(n \geqslant 1\) ;
\item
  The element
\end{enumerate}

\[
y _ {n} = \sum_ {k = 1} ^ {n} \gamma (1 - \gamma) ^ {k - 1} e _ {k} + (1 - \gamma) ^ {n} \cdot 1
\]

is invertible in \(\widetilde{\mathrm{A}}\) for all \(n\geqslant 1\) and

\[
\left\| y _ {n} ^ {- 1} x - y _ {n - 1} ^ {- 1} x \right\| <   \delta 2 ^ {- n} \quad (n \geqslant 2), \quad \left\| y _ {1} ^ {- 1} x - x \right\| <   \delta / 2.
\]

(Assuming that \((e_{1},\ldots,e_{n})\) have been defined, consider an element \(e_{n+1}\) arbitrary of norm \textless{} 2; observe that

\[
y _ {n + 1} = ((1 - \gamma) \cdot 1 + e _ {n + 1}) \left(\sum_ {k = 1} ^ {n} \gamma (1 - \gamma) ^ {k - 1} e _ {k} ^ {\prime} + (1 - \gamma) ^ {n} \cdot 1\right),
\]

where \(e_k' = ((1 - \gamma) \cdot 1 + e_{n+1})^{-1} e_k\) , and determine \(e_{n+1}\) satisfying the required conditions by exploiting the fact that the set of invertible elements of \(\widetilde{\mathsf{A}}\) is open, and that the map \(a \mapsto a^{-1}\) is continuous.)
\end{enumerate}

\begin{enumerate}
\def\labelenumi{\alph{enumi})}
\setcounter{enumi}{4}
\tightlist
\item
  For every \(x \in A\) , there exist y and z in A such that x = yz ("Cohen's factorization theorem"). Let \(\delta > 0\) . One may assume that z belongs to the closed left ideal generated by \(x\) and that \(\|y-z\|<\delta\) . (Consider a sequence \((e_{n})_{n\geqslant1}\) as above and define \(y_{n}\) as before; the sequence \(z_{n}=y_{n}^{-1}x\) converges to an element \(z\in\mathrm{A}\) , and \(z_{n}\) belongs to the closed left ideal generated by \(x\) ; the sequence \((y_{n})\) converges to \(y\in\mathrm{A}\) and \(x=yz\) .)
\end{enumerate}

\begin{enumerate}
\def\labelenumi{\arabic{enumi})}
\setcounter{enumi}{30}
\tightlist
\item
  Let A be the Banach space of complex integrable functions on {[}0, 1{]} with respect to Lebesgue measure. Set
\end{enumerate}

\[
(f \cdot g) (t) = \int_ {0} ^ {t} f (t - x) g (x) d x
\]

for \(t \in [0, 1]\) .

\begin{enumerate}
\def\labelenumi{\alph{enumi})}
\item
  The function \(f \cdot g\) is defined almost everywhere and is an element of A.
\item
  Equipped with the operation \((f,g)\mapsto f\cdot g\) , the space A is a commutative Banach algebra.
\item
  Let \(x \in \mathrm{A}\) be the constant function equal to 1. The function \(x^n\) is
\end{enumerate}

\[
t \mapsto \frac {t ^ {n - 1}}{(n - 1) !}.
\]

\begin{enumerate}
\def\labelenumi{\alph{enumi})}
\setcounter{enumi}{3}
\item
  The algebra A is topologically generated by x, the element x is quasi-nilpotent, and A is equal to its radical.
\item
  The algebra A admits an approximate identity.
\item
  The algebra A has no maximal ideal, regular or not, closed or not (use the preceding exercise and exercise 1 of I, p.~166). In particular, X(A) is empty, hence compact.
\end{enumerate}

\begin{enumerate}
\def\labelenumi{\arabic{enumi})}
\setcounter{enumi}{31}
\tightlist
\item
  Let G be a group. One says that G has the Powers property if, for every \(n \in N\) and every finite subset F of \(G = \{e\}\) , there exists a partition of G into two subsets A and B and a family \((g_i)_{1 \leqslant i \leqslant n}\) of elements of G such that
\end{enumerate}

\begin{enumerate}
\def\labelenumi{(\roman{enumi})}
\item
  for all \(g \in \mathrm{F}\) , the sets A and gA are disjoint;
\item
  for all distinct integers \(j\) and \(k\) between 1 and \(n\) , the sets \(g_{j}\mathrm{B}\) and \(g_{k}\mathrm{B}\) are disjoint.
\end{enumerate}

Let G be a group possessing the Powers property. Let I be a nonzero closed two-sided ideal of the reduced C*-algebra \(\text{Stell}_{r}(\text{G})\) (exercise 25). Let \(\pi: C[G] \to \text{Stell}_{r}(G)\) be the canonical representation.

\begin{enumerate}
\def\labelenumi{\alph{enumi})}
\tightlist
\item
  There exists a family \((\lambda_g)_{g\in \mathrm{G - }\{e\}}\) into \(\mathbf{C}\) such that
\end{enumerate}

\[
x = \sum_ {g \in \mathrm{G} - \{e \}} \lambda_ {g} \pi (g)
\]

is a hermitian element of \(\operatorname{Stell}_{r}(\mathrm{G})\) such that \(y = 1 + x\) belongs to I.

\begin{enumerate}
\def\labelenumi{\alph{enumi})}
\setcounter{enumi}{1}
\tightlist
\item
  Let \(\varepsilon > 0\) such that \(\varepsilon < 1\) . There exists a finite subset F of G = \{e\} and an integer \(n \geqslant 1\) such that the element
\end{enumerate}

\[
x ^ {\prime} = \sum_ {g \in \mathrm{F}} \lambda_ {g} \pi (g)
\]

satisfies \(\|x' - x\| < \varepsilon\) and \(\|x'\| < \frac{\sqrt{n}}{2}(1 - \varepsilon)\) .

\begin{enumerate}
\def\labelenumi{\alph{enumi})}
\setcounter{enumi}{2}
\tightlist
\item
  Let \(G = A \cup B\) and \((g_{i})_{1 \leqslant i \leqslant n}\) satisfying the conditions of the Powers property for the finite subset F and the integer n.~Set
\end{enumerate}

\[
y ^ {\prime} = \frac {1}{n} \sum_ {i = 1} ^ {n} \pi (g _ {i}) x ^ {\prime} \pi (g _ {i} ^ {- 1}).
\]

We have \(\|y'\| < 1 - \varepsilon\) . (Let \(p_i\) be the orthogonal projection of \(\ell^{2}(\mathrm{G})\) on \(\ell^{2}(g_{i}\mathrm{B})\) ; verify that

\[
(1 - p _ {i}) \pi (g _ {i}) x ^ {\prime} \pi (g _ {i} ^ {- 1}) (1 - p _ {i}) = 0
\]

for \(1 \leqslant i \leqslant n\) , and deduce that

\[
1 + y ^ {\prime} = 1 + \frac {1}{n} \sum_ {i = 1} ^ {n} p _ {i} \pi (g _ {i}) x ^ {\prime} \pi (g _ {i} ^ {- 1}) + \left(\frac {1}{n} \sum_ {i = 1} ^ {n} p _ {i} \pi (g _ {i}) x ^ {\prime} \pi (g _ {i} ^ {- 1}) (1 - p _ {i})\right) ^ {*},
\]

then note that the endomorphisms \(p_{i}\pi(g_{i})x^{\prime}\pi(g_{i}^{-1})\) have pairwise orthogonal images to estimate their norms.)

\begin{enumerate}
\def\labelenumi{\alph{enumi})}
\setcounter{enumi}{3}
\tightlist
\item
  The ideal I contains the element
\end{enumerate}

\[
1 + \frac {1}{n} \sum_ {i = 1} ^ {n} \pi (g _ {i}) x \pi (g _ {i} ^ {- 1}),
\]

and this element is invertible. Consequently, the stellar algebra \(\text{Stell}_{r}(\text{G})\) contains no closed two-sided ideal other than \(\{0\}\) and of \(\text{Stell}_{r}(\text{G})\) .

\begin{enumerate}
\def\labelenumi{\alph{enumi})}
\setcounter{enumi}{4}
\tightlist
\item
  Let G be a noncommutative free group. Then G has the Powers property.
\end{enumerate}
